\documentclass[11pt]{article}
% LaTeX-only set-up (the web build does not read this file).
\usepackage[T1]{fontenc}
\usepackage[utf8]{inputenc}
\usepackage{lmodern}
\usepackage{amsmath,amssymb,amsthm}
\usepackage[a4paper,margin=1.1in]{geometry}
\usepackage{microtype}
\usepackage{booktabs}
\usepackage{enumitem}
\usepackage{newunicodechar}
\usepackage[hidelinks]{hyperref}
\setlist{itemsep=0.2em,topsep=0.4em}
% Statements carry their own explicit label through \head, so nothing is numbered automatically.
\newenvironment{statement}{\par\addvspace{\medskipamount}\noindent\ignorespaces}{\par\addvspace{\medskipamount}}
\newenvironment{frontnote}{\par\addvspace{\medskipamount}\small\noindent\ignorespaces}{\par\addvspace{\medskipamount}}
\renewcommand{\qedsymbol}{$\blacksquare$}
% Unicode that may appear inside verbatim and \texttt (Lean code)
\newunicodechar{α}{\ensuremath{\alpha}}\newunicodechar{β}{\ensuremath{\beta}}
\newunicodechar{γ}{\ensuremath{\gamma}}\newunicodechar{δ}{\ensuremath{\delta}}
\newunicodechar{μ}{\ensuremath{\mu}}\newunicodechar{ω}{\ensuremath{\omega}}
\newunicodechar{τ}{\ensuremath{\tau}}\newunicodechar{λ}{\ensuremath{\lambda}}
\newunicodechar{Π}{\ensuremath{\Pi}}\newunicodechar{Σ}{\ensuremath{\Sigma}}
\newunicodechar{∀}{\ensuremath{\forall}}\newunicodechar{∃}{\ensuremath{\exists}}
\newunicodechar{∧}{\ensuremath{\wedge}}\newunicodechar{∨}{\ensuremath{\vee}}
\newunicodechar{¬}{\ensuremath{\neg}}\newunicodechar{→}{\ensuremath{\to}}
\newunicodechar{↔}{\ensuremath{\leftrightarrow}}\newunicodechar{≤}{\ensuremath{\le}}
\newunicodechar{∈}{\ensuremath{\in}}\newunicodechar{⊢}{\ensuremath{\vdash}}
\newunicodechar{₂}{\ensuremath{_2}}\newunicodechar{₃}{\ensuremath{_3}}\newunicodechar{⁺}{\ensuremath{^+}}
% Lean listings are wider than the text font allows, so set them smaller.
\makeatletter\def\verbatim@font{\normalfont\ttfamily\footnotesize}\makeatother

% Shared by the PDF and the web build. Keep to simple \newcommand definitions.
% The short name of the main notion (alpha-plus-recursively Mahlo). Change it here and nowhere else.
% It is used both as an adjective ("the least \RM{} ordinal") and as a noun at the start of
% sentences, so it has to be an abbreviation or symbol, not a lower-case word.
\newcommand{\RM}{RM\textsuperscript{+}}
% Bold run-in head of a statement or front note: \head{Lemma 4.1}
\newcommand{\head}[1]{\textbf{#1.}}
% Named classes and predicates
\newcommand{\Mh}{\mathrm{Mh}}
\newcommand{\Sat}{\mathrm{Sat}}
\newcommand{\Mem}{\mathrm{Mem}}

\title{Where the parameter-free $\delta$ first exceeds Girard's $\Xi_1$}
% The date is stated here once; the PDF and the web page both take it from this line.
\newcommand{\paperversion}{2026-11-22}
\date{\paperversion}
\begin{document}
\maketitle
\newcommand{\Xo}{\Xi_1}
\newcommand{\Xp}{\Xi'_1}
\newcommand{\PR}{\mathrm{PAD}_{\mathcal R}}
\newcommand{\PRp}{\mathrm{PAD}'_{\mathcal R}}
\newcommand{\dz}[1]{\delta^0_{#1}}
\newcommand{\nx}{\operatorname{nx}}
\newcommand{\Pre}{\mathrm{Pre}}
\newcommand{\Dil}{\mathrm{Dil}}
\newcommand{\tl}{\triangleleft}
\newcommand{\E}{\mathrel{E}}

\begin{frontnote}\head{Abstract} For an infinite ordinal $\alpha$ let $\Xo(\alpha)$ be the supremum of the ordinals
$D(\alpha)$ over all recursive dilators $D$; it is the least ordinal above $\alpha$ that is closed under every
recursive dilator. Let $\Xp(\alpha)$ be the same supremum over all recursive predilators that are well-ordered at
$\alpha$. The notation $\Xi_1$ is Girard's; Ressayre (1982) writes $\mathrm{PAD}$ and $\mathrm{PAD}'$.
$\Xp(\alpha)$ is also the supremum of the well-orderings
that are $\Sigma_1$-definable over $L_{\alpha}$ without parameters. Ressayre showed that
$\Xo(\alpha) < \Xp(\alpha)$ can happen. We determine where it first happens. Let $u$ be the least ordinal that is
uncountable in the next admissible set $L_{u^+}$, and let $\mu_0$ be the least ordinal that is not
$\Sigma_1$-definable from $u$ in $L_{u^+}$. Then $\Xp(d) = \Xo(d)$ for every infinite $d < \mu_0$, while
$\Xo(x) < \Xp(x) = \Xo(\Xo(x)) < \mu_0^+$ for every $x$ with $\mu_0 \le x < \Xo(\mu_0)$. We give a criterion for
$\Xp(\alpha) = \Xo(\alpha)$ at an arbitrary $\alpha$, and characterise $\mu_0$ as the least ordinal that shares all
its $\Sigma^1_1$ properties with a larger ordinal. As an application we show that a $\Sigma^1_1$ condition on a
family of partial operations makes the least ordinal closed under the family equal to $\Xo(\alpha)$, and we record,
partly from cited results, where the next admissible ordinal, the supremum $\delta_{\alpha}$ of the well-orderings $\Delta_1$-definable over
$L_{\alpha}$ with parameters, and the parameter-free $\Sigma_1$ form stop being equal to $\Xo(\alpha)$.
\end{frontnote}

\begin{frontnote}\head{Provenance} The proofs were found and written by an AI system. No human specialist has
reviewed them, and novelty is not established. Section 12 states what has been checked and how.
\end{frontnote}

\section{Introduction}

Girard's dilators are the functors on the category of linear orders that preserve direct limits and pullbacks and
send well-orders to well-orders; a predilator need not send well-orders to well-orders (Section 2.1). For an
infinite ordinal $\alpha$ put
\begin{align*}
\Xo(\alpha) &= \sup\{|D(\alpha)| : D \text{ a recursive dilator}\}, \\
\Xp(\alpha) &= \sup\{|P(\alpha)| : P \text{ a recursive predilator, } P(\alpha) \text{ well-ordered}\},
\end{align*}
where $|W|$ is the order type of a well-order $W$. The notation is Girard's: in Volume 2 of his book (Corollary 11.B.7) $\Xi_1$ is the sum of all
recursive dilators, whose value at $\alpha$ is $\Xo(\alpha)$, and $\Xi'_1$ follows his Exercise 11.B.12. Ressayre
(1982) writes $\mathrm{PAD}$ and $\mathrm{PAD}'$. He defines the two ordinals
with a slightly different class of functors, the primitive recursive functors on the finite ordinals extended by
direct limits; we write $\PR(\alpha)$ and $\PRp(\alpha)$ for his ordinals, and Corollary 5 shows that
$\PR(\alpha) = \Xo(\alpha)$ and $\PRp(\alpha) = \Xp(\alpha)$ for infinite $\alpha$. At a finite argument $n$ both
of Ressayre's ordinals are $\omega$, because his functors have finite values at $n$ and the constant functors are
among them; so only infinite arguments are considered.

Throughout, $\alpha^+$ is the least admissible ordinal above $\alpha$. Girard proved that
$\Xo(\alpha) = \alpha^+$ for many ordinals $\alpha$, among them every infinite $\alpha$ below the ordinal $\sigma$
defined in Section 1.1 (Girard's book, Volume 2, Theorem 11.B.6 and Corollary 11.B.7). Ressayre gave a new proof and asked how $\Xo(\alpha)$, $\Xp(\alpha)$ and $\alpha^+$ can be
related. He showed $\Xo(\alpha) \le \Xp(\alpha) \le \alpha^+$, proved equality of the first two in several cases,
and gave the first example of inequality (his Theorem IV.7): if $\lambda$ is the least ordinal that is not
countable in $L_{\lambda^+}$, and $\mu$ is the least ordinal above $\lambda$ that is not $\Sigma_1$-definable from
$\lambda$ in $L_{\lambda^+}$, then
\[ \PR(\lambda) = \PR(\mu) < \PRp(\mu) < \mu^+ . \]
In the proof he writes $\mu = \lambda + p_0$, where $p_0 < \lambda$ is the least ordinal that is not
$\Sigma_1$-definable from $\lambda$ in $L_{\lambda^+}$; he does not evaluate the two functions at $p_0$.

This paper determines the least ordinal at which the two functions differ. It is $p_0$, which we call $\mu_0$.

\subsection{Results}

\emph{Definability.} A set is \emph{$\Sigma_1$-definable from $p$ in $L_{\rho}$} if it is the only element $x$ of
$L_{\rho}$ with $L_{\rho} \models \psi(x, p)$, for some $\Sigma_1$ formula $\psi$. A \emph{$\Sigma^1_1$ class} is a
class of ordinals of the form $\{\beta : (\beta, <) \models \exists \vec R\, \theta(\vec R)\}$ with $\theta$
first-order. Following Ressayre, an ordinal is \emph{$\beta^-$-definable} if it is the least element of a
$\Sigma^1_1$ class, and \emph{$\beta$-definable} if it is the only element of one. An infinite ordinal is
$\beta^-$-definable exactly when it is the least ordinal at which some recursive predilator is not well-ordered
(Lemma 3.4).

\emph{Four ordinals.} All four are countable.
\begin{itemize}
\item $\boldsymbol{\sigma}$ is the least ordinal that is not $\beta$-definable. By Ressayre (III.10) it is the least
  $\Sigma^1_1$-reflecting ordinal; we use only the definition just given.
\item $\boldsymbol{u}$ is the least ordinal $\alpha$ that is uncountable in $L_{\alpha^+}$. This is Ressayre's
  $\lambda$; Section 2.4 mentions related ordinals from the literature.
\item $\boldsymbol{\mu_0}$ is the least ordinal that is not $\Sigma_1$-definable from $u$ in $L_{u^+}$. This is
  Ressayre's $p_0$. It satisfies $\mu_0 < u$, and it is the image of $u$ when the set $H$ of all elements of
  $L_{u^+}$ that are $\Sigma_1$-definable from $u$ is collapsed to a transitive set.
\item $\boldsymbol{\gamma}$ is the order type of the set of ordinals in $H$. The transitive collapse of $H$ is
  $L_{\gamma}$, and $\gamma$ is the least ordinal such that $\mu_0$ is countable in $L_{\gamma+1}$ (Lemma 2.6).
\end{itemize}
They are ordered as follows (Lemmas 2.5 and 2.6, Theorem 1, Proposition 6):
\[ \omega < \sigma < \sigma + \omega < \mu_0 < \gamma = \Xo(\mu_0) < \Xo(\gamma) < \mu_0^+ < u < u^+ . \]

\emph{Closure points and definable well-orderings.} An ordinal $\kappa$ is a \emph{closure point of the recursive dilators}
if $|D(\beta)| < \kappa$ for every $\beta < \kappa$ and every recursive dilator $D$; $\Xo(\alpha)$ is the least
closure point above $\alpha$ (Lemma 2.1). For an infinite ordinal $\alpha$:
\begin{itemize}
\item $\delta_{\alpha}$ is the supremum of the order types of the well-orderings of subsets of $L_{\alpha}$ that
  are $\Delta_1$-definable over $L_{\alpha}$ with parameters from $L_{\alpha}$. The definition is modelled on Aguilera's
  $\delta_{\alpha}$ for admissible $\alpha$ (The order of reflection, Section 2), the supremum of the order types of
  the $\alpha$-recursive well-orderings of subsets of $\alpha$; there $\alpha$ is called Gandy when
  $\delta_{\alpha} = \alpha^+$. We use the definition above at every infinite $\alpha$, admissible or not, and no
  result about Aguilera's $\delta_{\alpha}$.
\item $\dz{\alpha}$ is the supremum of the order types of the well-orderings of subsets of $L_{\alpha}$ that are
  $\Sigma_1$-definable over $L_{\alpha}$ without parameters.
\end{itemize}
The two definitions differ in two respects: $\dz{\alpha}$ allows no parameters, and it allows $\Sigma_1$ in place of
$\Delta_1$ definitions. The first difference is the one that matters in this paper. The second is made because the
$\Sigma_1$ form is the one we can identify with $\Xp(\alpha)$ at every infinite $\alpha$; at limit $\alpha$ the
parameter-free $\Delta_1$ form gives the same ordinal (remark after Lemma 6.1).

\begin{statement}\head{Theorem 1} (i) $\Xp(d) = \Xo(d)$ for every infinite ordinal $d < \mu_0$.

(ii) $\mu_0$ is a closure point of the recursive dilators, $\Xo(\mu_0) = \gamma$, and for every $x$ with
$\mu_0 \le x < \gamma$,
\[ \Xo(x) = \gamma < \Xp(x) = \Xo(\gamma) < \mu_0^+ . \]

(iii) $\Xp(\gamma) = \Xo(\gamma)$.
\end{statement}

So $\mu_0$ is the least infinite ordinal at which the two functions differ, they differ on the whole interval
$[\mu_0, \gamma)$, and there $\Xp(x) = \Xo(\Xo(x))$: exactly one closure point, $\gamma$, lies between $x$ and
$\Xp(x)$. At $\gamma$ itself the functions agree again.

\begin{statement}\head{Theorem 2} (criterion) Let $\alpha$ be an infinite ordinal.

(a) Let $F$ be a recursive predilator with $F(\alpha)$ well-ordered. Then $|F(\alpha)| < \Xo(\alpha)$ if and only if
there is a $\Sigma^1_1$ class $X$ with $\alpha \in X$ such that $F(\beta)$ is well-ordered for every $\beta \in X$.

(b) $\Xp(\alpha) = \Xo(\alpha)$ if and only if either no $\beta^-$-definable ordinal is greater than $\alpha$, or
some $\Sigma^1_1$ class containing $\alpha$ has all its elements below the least $\beta^-$-definable ordinal greater
than $\alpha$.
\end{statement}

\begin{statement}\head{Theorem 3} (characterisations of $\mu_0$) (a) Every $\Sigma^1_1$ class that contains $\mu_0$
contains $u$.

(b) $\mu_0$ is the least ordinal $d$ such that every $\Sigma^1_1$ class containing $d$ has an element $\ge u$.

(c) $\mu_0$ is the least ordinal $d$ for which there is an ordinal $y > d$ that belongs to every $\Sigma^1_1$ class
containing $d$.
\end{statement}

Theorem 3 (a) is the reason for the failure at $\mu_0$: a predilator that first fails at $u$ is well-ordered at
$\mu_0$, but by Theorem 2 (a) no recursive dilator bounds it there.

\subsection{The next closure point}

Girard's theorem identifies the next admissible ordinal $\alpha^+$ with the least closure point of the recursive
dilators above $\alpha$, for $\alpha < \sigma$. Ordinal collapsing functions are built over a sequence of such
``next'' ordinals, each closed under the operations available below it, and the question behind this paper is
which function of $\alpha$ gives the least closure point above $\alpha$ in general. Besides $\Xo(\alpha)$ itself,
three functions are candidates: $\alpha^+$, the ordinal $\delta_{\alpha}$ modelled on Aguilera's $\delta_{\alpha}$, and its
parameter-free form $\dz{\alpha} = \Xp(\alpha)$. This paper does not define a collapsing function. It records where
each candidate stops being the least closure point (table below), and first gives a result that explains the outcome.

Suppose that for each infinite $\beta$ a set $\mathcal O_\beta$ of recursive predilators
well-ordered at $\beta$ is given, the operations \emph{permitted at} $\beta$. Let $f_{\mathcal O}(\alpha)$ be the
least $\kappa > \alpha$ such that $|P(\beta)| < \kappa$ whenever $\omega \le \beta < \kappa$ and
$P \in \mathcal O_\beta$.

\begin{statement}\head{Theorem 4} (a) Suppose that every recursive dilator belongs to every $\mathcal O_\beta$, and
that every $P \in \mathcal O_\beta$ is well-ordered at every element of some $\Sigma^1_1$ class containing $\beta$.
Then $f_{\mathcal O}(\alpha) = \Xo(\alpha)$ for every infinite $\alpha$.

(b) Let $\mathcal O^*_\beta$ be the set of all recursive predilators that are well-ordered at $\beta$. Then
$\mathcal O^*$ satisfies the first hypothesis of (a), and it satisfies the second with ``$\Pi^1_1$ class'' in
place of ``$\Sigma^1_1$ class''. For it, $f_{\mathcal O^*}(\alpha) = \Xo(\alpha)$ for all infinite
$\alpha < \mu_0$, and $f_{\mathcal O^*}(\mu_0) > \Xo(\mu_0)$.
\end{statement}

Part (a) says that $\Xo$ results whenever each partial operation is permitted at $\beta$ only together with a
$\Sigma^1_1$ property of $\beta$ that guarantees it is well-ordered. Part (b) says that the $\Sigma^1_1$ condition
cannot be dropped or replaced by a $\Pi^1_1$ one; $f_{\mathcal O^*}(\alpha)$ is the least ordinal above $\alpha$
closed under $\beta \mapsto \dz{\beta}$.

The other two results compare the candidates with $\Xo$.

\begin{statement}\head{Corollary 5} For every infinite ordinal $\alpha$,
\[ \Xo(\alpha) = \PR(\alpha), \qquad \Xp(\alpha) = \PRp(\alpha) = \dz{\alpha} . \]
Hence the least infinite $\alpha$ with $\PR(\alpha) \ne \PRp(\alpha)$ is Ressayre's $p_0$, with
$\PRp(p_0) = \PR(\PR(p_0)) < p_0^+$; and the least infinite $\alpha$ at which
$\dz{\alpha}$ is not the least closure point of the recursive dilators above $\alpha$ is $\mu_0$.
\end{statement}

\begin{statement}\head{Proposition 6} Let $\sigma$ be the least ordinal that is not $\beta$-definable. Then
$\sigma < \mu_0$ and
\[ \dz{\sigma + \omega} = \Xo(\sigma) < \sigma^+ = \delta_{\sigma + \omega} , \]
while $\sigma < \sigma + \omega < \delta_{\sigma} < \sigma^+$. In particular $\delta_{\alpha}$ is not constant on
$[\sigma, \delta_{\sigma})$.
\end{statement}

Since $\Xo(\sigma + \omega) = \Xo(\sigma)$ by Lemma 2.1, the proposition gives
$\delta_{\sigma + \omega} > \Xo(\sigma + \omega)$: with parameters, $\delta_{\alpha}$ exceeds the least closure point
already at $\sigma + \omega$. A function that sends $\alpha$ to the least element above $\alpha$ of a fixed class of
ordinals is constant on $[\alpha, f(\alpha))$, so $\alpha \mapsto \delta_{\alpha}$ is not such a function either.
\begin{center}\begin{tabular}{p{0.13\linewidth}p{0.42\linewidth}p{0.33\linewidth}}
\toprule
Candidate & First infinite $\alpha$ where it is not $\Xo(\alpha)$ & Basis \\
\midrule
$\alpha^+$ & $\sigma$, where $\Xo(\sigma) < \sigma^+$ & cited (Ressayre, I.6 and IV.6) \\
$\delta_{\alpha}$ & at most $\sigma + \omega$ & Proposition 6 (cited results and a construction) \\
$\dz{\alpha}$ & $\mu_0$, where $\dz{\mu_0} = \Xo(\Xo(\mu_0))$ & Theorem 1, Corollary 5 \\
\bottomrule
\end{tabular}\end{center}

\subsection{Overview of the proofs}

\emph{The criterion.} A $\Sigma^1_1$ condition on a well-order $b$ is the existence of an infinite branch through a
tree whose nodes carry finitely many labels from $b$ and whose membership test reads only the order pattern of the
labels. The nodes of such a tree, in the Kleene--Brouwer order (Section 3.1), form a recursive predilator in $b$
(Lemma 3.1). If the nodes also follow a relation $\tl$ that the guessed structure defines, the rank of the tree is
at least the height of $\tl$, that is, its rank as a well-founded relation (Lemma 3.2). For Theorem 2 (a), right to
left, the guessed structure is a witness for $\beta \in X$ and $\tl$ is the order of $F(\beta)$: the tree is
well-founded at every ordinal because $F$ is well-ordered on $X$, so it gives a recursive dilator, and its value
at $\alpha$ exceeds $|F(\alpha)|$. This is the argument of Ressayre's Theorem IV.6, which he credits in its
essential step to Harrington, with a weaker hypothesis and with trees in place of models of set theory. Left to
right is immediate: if a dilator $D$ bounds $F$ at $\alpha$, the class of $\beta$ at which $F(\beta)$ embeds in
$D(\beta)$ is $\Sigma^1_1$.

\emph{Failure at $\mu_0$.} The collapse of $H$ sends $u$ to $\mu_0$ and every $\Sigma_1$ fact about $u$ in
$L_{u^+}$ to a true fact about $\mu_0$. Since $\Sigma^1_1$ properties of a countable ordinal $\beta$ are $\Pi_1$
properties of $\beta$ in $L_{\beta^+}$, every $\Sigma^1_1$ class containing $\mu_0$ contains $u$ (Lemma 8.3). So
by Theorem 2 (a) a recursive predilator $F_0$ that first fails at $u$ has $|F_0(\mu_0)| \ge \Xo(\mu_0)$.

\emph{Agreement below $\mu_0$.} Let $d < \mu_0$ and let $r$ be the least $\beta^-$-definable ordinal above $d$; then
$r \le u$. The set $Z_d$ of indices of recursive predilators whose value at $d$ is a well-order of type below
$\Xo(d)$ contains every index of a recursive dilator, and it is $\Sigma_1$-definable from $r$ in $L_{r^+}$. A
self-referential sentence then produces an index in $Z_d$ of a predilator that is well-ordered exactly below $r$
(Lemma 7.3). By Theorem 2 (a) that predilator is well-ordered on a $\Sigma^1_1$ class containing $d$, which
therefore lies below $r$, and Theorem 2 (b) gives equality. The lemma needs $r$ to be $\beta$-definable; for
$r < u$ this is Lemma 7.1, and for $r = u$ it is a remark of Ressayre (Lemma 7.2).

\emph{The values.} $\Xo(\mu_0) = \gamma$ because the isomorphisms between the orders $D(u)$ and their order types
lie in $H$ and collapse to isomorphisms for $D(\mu_0)$ (Lemma 8.4). On $[\mu_0, u)$ the predilators that are
well-ordered do not depend on the argument, since no $\beta^-$-definable ordinal lies in that interval (Lemma 8.6).
Composition with $F_0$ gives $\Xp(\mu_0) \ge \Xp(\gamma)$, an explicit $\Sigma^1_1$ class containing $\gamma$ and
bounded by $u$ gives $\Xp(\gamma) = \Xo(\gamma)$, and a collection argument gives $\Xo(\gamma) < \mu_0^+$
(Lemma 8.5).

\begin{center}\begin{tabular}{p{0.30\linewidth}p{0.64\linewidth}}
\toprule
Result & Uses \\
\midrule
Lemmas 2.1--2.6 & Facts (L1)--(L3); Ressayre, Remark IV.5 \\
Lemmas 3.1--3.4 & Sections 2.1 and 2.2 \\
Theorem 2 & Lemmas 3.1--3.4 \\
Lemmas 5.1, 5.2, Corollary 5.3 & Lemmas 2.2, 3.1--3.3; (L1), (L3), (L4) \\
Lemmas 6.1, B.1, B.2 & Lemmas 3.1--3.3 \\
Corollary 5 & Lemmas 2.1, 6.1, B.2; Theorem 1 for the last sentence \\
Lemma 7.1 & Lemmas 2.3, 3.3 (C), 3.4, 5.1; (L3) \\
Lemma 7.2 & Lemma 2.4, Corollary 5.3 (d); (L1) \\
Lemma 7.3 & Lemmas 3.1, 3.3 (C), Corollary 5.3 (d) \\
Lemmas 8.1--8.6 & Theorem 2; Lemmas 2.4--2.6, 3.4, 7.2, 7.3; Corollary 5.3; (L3), (L5) \\
Theorem 1 & Theorem 2; Lemmas 2.1, 2.3--2.6, 3.3 (L), 7.1--7.3, 8.1--8.6 \\
Theorem 3 & Lemmas 2.3--2.6, 8.2, 8.3; Corollary 5.3 (d) \\
Theorem 4 & Theorems 1 and 2; Lemmas 2.1, 3.4 \\
Proposition 6 & Theorem 1, Corollary 5, Lemmas 2.1, 2.5; (L2), (L5); Ressayre I.6, III.6 (a), III.10, the example
after IV.6; Sacks \\
\bottomrule
\end{tabular}\end{center}

\subsection{Relation to the literature}

\emph{Ressayre (1982).} The ordinals $u$ and $\mu_0$ are defined there; he shows that $u$ is $\beta$-definable,
that $\mu_0 < u$, and Theorem IV.7 as quoted above (Remark IV.5, Theorem IV.7). He also shows that the two
functions agree at every $\beta$-definable ordinal, at every ordinal that is the last element of a $\Sigma^1_1$
class, and at every ordinal above the least stable ordinal (the least $\alpha$ such that $L_{\alpha}$ is a
$\Sigma_1$-elementary substructure of $L$), and that they differ cofinally often below it (IV.3, IV.6, IV.8). We
did not find in that paper the values of the two functions at $p_0$, the question of the least ordinal at which
they differ, Theorem 2 (b) as a criterion, or Theorem 3. Theorem 2 (a), right to left, is the proof of his IV.6
with the hypothesis ``$\alpha$ is the last element of a $\Sigma^1_1$ class'' replaced by ``$F$ is well-ordered on a
$\Sigma^1_1$ class containing $\alpha$''. The self-referential sentence in Lemma 7.3 follows the method of his
proof of Girard's theorem, which he attributes to Harrington.

\emph{Other sources.} Gostanian, and Aguilera in a sharper form, determine when an admissible
$\alpha$ is Gandy, that is, when Aguilera's $\delta_{\alpha}$ equals $\alpha^+$; we do not use their results.
Gostanian calls such $\alpha$ good. His method for proving an ordinal good (his Proposition 2.1) lifts Gandy's
theorem on recursive linear orderings that are not well-ordered but have no hyperarithmetic descending sequence
(Gandy 1960) from $\omega$ to every admissible $\alpha$. The name Gandy, introduced by Abramson and Sacks (1976),
refers to that theorem. Aguilera and Pakhomov use recursive predilators to relate parameter-free
$\Sigma^1_1$-reflection to the $\Pi^1_2$ consequences of theories; we know of no comparison there between
$\dz{\alpha}$ and closure points of dilators. Exercise 11.B.12 of Volume 2 of Girard's book compares
$\Xi_1(\sigma)$ with the corresponding sum $\Xi'_1(\sigma)$ over the predilators well-ordered at $\sigma$, quoting
Ressayre for their equality. Girard calls the ordinals
of Lemma 3.4 weakly $\beta$-definable (Remark 11.B.5). We found no comparison of the two functions above $\sigma$
in Annexes 11.B and 11.C. We have not read Girard and Ressayre (1985), the rest of Girard's book, or the later
literature on dilators with this question in mind.
Novelty of Theorems 1, 3 and 4 is not established.

\section{Preliminaries}

Ordinals are denoted by $\alpha, \beta, \eta, \xi, \rho, \zeta, \kappa$; $d$, $x$ and $y$ are also used for ordinals
where stated. In Ressayre's terms ``$\beta$-definable'' and ``$\beta^-$-definable'' the letter $\beta$ is part of the
name and is not a variable. \emph{Countable} means finite or countably infinite.

\subsection{Recursive predilators and dilators}

We use Girard's normal form (Girard 1981, Theorem 2.3.12 and Remarks 2.3.13). The \emph{order pattern} of two finite tuples $s, s'$ from a
linear order is the total preorder that the order induces on their positions; for given lengths there are finitely
many order patterns. A \emph{recursive predilator} $P$ is given by a recursive set $\mathrm{Tr} \subseteq \omega$ of
\emph{traces}, a recursive function $\mathrm{ar} : \mathrm{Tr} \to \omega$, the \emph{arity}, and a recursive
\emph{comparison}, which assigns one of $<$, $=$, $>$ to two traces $z, z'$ together with an order pattern of an
increasing $\mathrm{ar}(z)$-tuple and an increasing $\mathrm{ar}(z')$-tuple. For a linear order $b$, $P(b)$ is the
set of pairs $(z, s)$ with $z \in \mathrm{Tr}$ and $s$ an increasing $\mathrm{ar}(z)$-tuple from $b$, compared by the
comparison applied to $z$, $z'$ and the order pattern of $s$ and $s'$. It is required that $P(b)$ is a strict linear
order for every $b$, and that $P(f)(w) \le P(g)(w)$ for all $w \in P(b)$ whenever $f, g : b \to b'$ are embeddings
with $f \le g$ pointwise, where $P(f)(z, s) = (z, f \circ s)$. The second requirement is part of Girard's
definition of a predilator (Girard 1981, Definition 4.4.1(iii)); it is verified in Lemma 3.1 so that the functors built there are predilators in his sense, and it is
not used otherwise. A recursive predilator is a functor on linear orders that preserves direct limits and
pullbacks, and $P(n)$ may be infinite for finite $n$. A \emph{recursive dilator} is a recursive predilator $D$ such
that $D(b)$ is a well-order for every well-order $b$. Sums and compositions of recursive predilators are recursive
predilators, with indices computed primitive recursively.

For example, let $P(b)$ be a copy of the negative integers followed by $b$: the traces are one trace of arity $0$
for each negative integer and one trace of arity $1$. This is a recursive predilator and $P(b)$ is never a
well-order. Lemma 3.4 gives recursive predilators that are well-ordered at some ordinals and not at others.

$\Pre$ is the set of all numbers $e$ that code a recursive predilator $P_e$ through indices for its three
components, and $\Dil$ is the set of all $e \in \Pre$ such that $P_e$ is a dilator. Neither set selects one index
for each functor. The set $\Pre$ is arithmetical, because $P(b)$ is a linear order for every $b$ as soon as it
is one for every finite $b$; so $\Pre \in L_{\omega + \omega}$.

If $b \subseteq b'$ then $P(b)$ is a suborder of $P(b')$. So for a recursive predilator $P$ the ordinals $b$ with
$P(b)$ well-ordered form an initial segment of the ordinals, and $|P(\beta)| \le |P(\beta')|$ when
$\beta \le \beta'$ are in it. We say that $P$ \emph{first fails at} $r$ if $r$ is the least ordinal with $P(r)$
not well-ordered. Since $P(n)$ can be infinite and even ill-founded for finite $n$, some statements below speak of
the least \emph{infinite} ordinal with a property; this is deliberate.

For an infinite ordinal $\alpha$, $\Xo(\alpha)$ and $\Xp(\alpha)$ are defined at the start of Section 1. Neither supremum
is attained, since $P + 1$ is a recursive predilator. Every recursive dilator is a recursive predilator
well-ordered at $\alpha$, so $\Xo(\alpha) \le \Xp(\alpha)$.

\begin{statement}\head{Lemma 2.1} $\Xo(\alpha)$ is the least closure point of the recursive dilators above $\alpha$.
Hence $\Xo$ is nondecreasing, and $\Xo(\beta) = \Xo(\alpha)$ whenever $\alpha \le \beta < \Xo(\alpha)$.
\end{statement}

\begin{proof} $\Xo(\alpha) > \alpha$ by the identity dilator. Let $\beta < \Xo(\alpha)$ and $D_2 \in \Dil$. Pick
$D_1 \in \Dil$ with $\beta < |D_1(\alpha)|$. Then $|D_2(\beta)| \le |(D_2 \circ D_1)(\alpha)| < \Xo(\alpha)$. So
$\Xo(\alpha)$ is a closure point. If $\kappa > \alpha$ is a closure point then $|D(\alpha)| < \kappa$ for every
$D \in \Dil$, so $\Xo(\alpha) \le \kappa$. The rest follows.
\end{proof}

\subsection{$\Sigma^1_1$ classes}

A $\Sigma^1_1$ class is defined in Section 1.1; we say that the first-order sentence $\theta$ \emph{defines} it, and
that an \emph{expansion} of a linear order $A$ is a structure $(A, <, \vec R)$. The language of $\theta$ is finite.
Constants and function symbols may occur in $\theta$; they stand for relation symbols together with the axioms saying
that these are graphs of functions. The empty order counts as a structure, so $\{0\}$ is the $\Sigma^1_1$ class
defined by $\forall x\, (x \ne x)$, and every finite ordinal is the only element of a $\Sigma^1_1$ class. An order of type $\beta$ has an
expansion satisfying $\theta$ if and only if $\beta$ is in the class. $\Sigma^1_1$ classes are closed under finite
unions and intersections, and the class of all infinite ordinals is $\Sigma^1_1$. A \emph{$\Pi^1_1$ class} is the
complement of a $\Sigma^1_1$ class. The terms $\beta$-definable and $\beta^-$-definable are Ressayre's (1982,
pp.\ 253 and 266). For an infinite ordinal $\alpha$, $r(\alpha)$ denotes the least $\beta^-$-definable ordinal
greater than $\alpha$, if there is one.

\subsection{The constructible hierarchy}

An ordinal $\rho > \omega$ is admissible if $L_{\rho}$ is a model of Kripke--Platek set theory KP, and $<_L$ is the
canonical well-ordering of $L$. Kripke--Platek set theory and admissible sets are treated in Barwise (1975).
$\Sigma_1$-definability from a parameter is defined in Section 1.1. We use the following facts.
\begin{itemize}
\item[(L1)] There are a $\Sigma_1$ formula and a $\Pi_1$ formula that both define the relation $z = L_a$ in every
  admissible $L_{\rho}$, and likewise for $<_L$. If $y \in L_a$ and $x <_L y$ then $x \in L_a$. For a limit $\delta$,
  the elements of $L_{\delta}$ that are $\Sigma_1$-definable from $p$ form a $\Sigma_1$-elementary substructure, and
  every $\Sigma_1$-elementary substructure of $L_{\delta}$ is isomorphic to exactly one $L_{\gamma}$ (Barwise,
  Sections II.5 and II.6; Devlin, Chapter II).
\item[(L2)] (Kripke, Platek) Every element of $L_{\alpha^+}$ is $\Sigma_1$-definable in $L_{\alpha^+}$ from finitely
  many ordinals $\le \alpha$ (Ressayre, III.2 and III.3, who cites Barwise).
\item[(L3)] If a well-order belongs to an admissible set, so do its order type and the isomorphism onto it
  (Barwise, Chapter I, by $\Sigma$-recursion). Hence $\Xp(\alpha) \le \alpha^+$, since $P(\alpha) \in L_{\alpha^+}$.
\item[(L4)] (Truncation lemma) The well-founded part of a model of KP is admissible or is the whole model
  (Barwise, II.8.4).
\item[(L5)] For every $\Sigma^1_1$ class $X$ there is a $\Pi_1$ formula $\varphi$ such that for every countable
  infinite $\beta$: $\beta \in X$ if and only if $L_{\beta^+} \models \varphi(\beta)$ (Ressayre, Lemma III.5, by the
  Barwise completeness theorem; compare Barwise, Gandy and Moschovakis, Theorem 3.1, and the uniform form stated
  by Gostanian, p. 189).
\end{itemize}

A \emph{term} $t$ is given by a $\Delta_0$ formula $\tau_t(w, v, z)$, written $\tau$ when $t$ is fixed. For an
infinite ordinal $\beta$, $t(\beta)$ \emph{is defined}
if some pair $(w, v) \in L_{\beta^+}$ satisfies $\tau(w, v, \beta)$, and then $t(\beta)$ is the second component of
the $<_L$-least such pair. By (L1) the sets $\Sigma_1$-definable from $\beta$ in $L_{\beta^+}$ are exactly the
values $t(\beta)$ (Ressayre, III.1). In a model $M$ of KP + $V = L$ we write $t^M(z)$ for the second component of
the $<_L^M$-least pair of $M$ that satisfies $\tau(w, v, z)$ in $M$.

\begin{statement}\head{Lemma 2.2} For every $P \in \Pre$ there is a term $t_P$ such that for every infinite
$\beta$: $t_P(\beta)$ is defined if and only if $P(\beta)$ is a well-order, and then $t_P(\beta) = |P(\beta)|$.
\end{statement}

\begin{proof} $P(\beta) \in L_{\beta^+}$, and $P(\beta)$ is a well-order if and only if $L_{\beta^+}$ contains an
isomorphism of it onto an ordinal, by (L3). Let $\tau(w, v, z)$ say that $w$ is such an isomorphism of $P(z)$ onto
the ordinal $v$, with $P(z)$ computed inside $w$.
\end{proof}

\begin{statement}\head{Lemma 2.3} An infinite ordinal $\alpha$ is countable in $L_{\alpha^+}$ if and only if every
ordinal below $\alpha$ is $\Sigma_1$-definable from $\alpha$ in $L_{\alpha^+}$. In that case every element of
$L_{\alpha^+}$ is $\Sigma_1$-definable from $\alpha$ in $L_{\alpha^+}$.
\end{statement}

\begin{proof} If $\alpha$ is countable in $L_{\alpha^+}$, the $<_L$-least surjection $f : \omega \to \alpha$ is
$\Sigma_1$-definable from $\alpha$, and so is each $f(n)$; the last sentence then follows from (L2). Conversely,
suppose every $\xi < \alpha$ is $\Sigma_1$-definable from $\alpha$, so every $\xi < \alpha$ is $t(\alpha)$ for a
term $t$. For a term $t$ and $a < \alpha^+$ let $t_a(\alpha)$ be the second component of the $<_L$-least pair
$(w, v) \in L_a$ with $\tau_t(w, v, \alpha)$, if there is one. By (L1), if $t_a(\alpha)$ exists then
$t(\alpha) = t_a(\alpha)$, because every pair below a pair in $L_a$ is in $L_a$. The function that sends
$\xi < \alpha$ to the least $a$ such that $t_a(\alpha) = \xi$ for some term $t$ is total on $\alpha$ and $\Sigma_1$
over $L_{\alpha^+}$, so by $\Sigma_1$-collection it is bounded by some $a^* < \alpha^+$. Enlarge $a^*$ so that
$\alpha < a^* < \alpha^+$; least pairs found at earlier levels keep their values by (L1). Then
$t \mapsto t_{a^*}(\alpha)$, restricted to the terms for which this is an ordinal below $\alpha$, maps a set of
formulas onto $\alpha$. Its graph is an element of $L_{\alpha^+}$: with $L_{a^*}$ as a parameter, satisfaction of
$\Delta_0$ formulas in $L_{a^*}$ and the search in $L_{a^*}$ for a least pair are $\Delta_1$ over $L_{\alpha^+}$, so
$\Delta_1$-separation applies.
\end{proof}

\subsection{The ordinals $u$, $\mu_0$ and $\gamma$}

Let $u$ be the least ordinal $\alpha$ that is uncountable in $L_{\alpha^+}$. Such ordinals exist (the first
uncountable cardinal is one), and the least one is countable by the L\"owenheim--Skolem theorem.

\begin{statement}\head{Lemma 2.4} (a) If $\beta < u$ is infinite then $\beta$ is countable in $L_{\beta^+}$ and
$\beta^+ < u$ (Ressayre, Remark IV.5). So $u$ is a limit of admissible ordinals and every element of $L_u$ is
countable in $L_u$.

(b) $u$ is admissible, and no ordinal $\beta$ with $u < \beta < u^+$ is admissible.

(c) $|F(x)| < u$ for every $x < u$ and every recursive predilator $F$ with $F(x)$ well-ordered. So $u$ is a
closure point of the recursive dilators.
\end{statement}

\begin{proof} (a) The first statement is the minimality of $u$, and the second is proved in the place cited. (b)
$L_u$ satisfies all axioms of KP except possibly $\Delta_0$-collection, as $u$ is a limit of admissibles. If
collection failed, there would be $a \in L_u$ and a map $f$ from $a$ cofinally into $u$ that is definable over
$L_u$, so $f \in L_{u+1}$. Then $a$ and each $f(z)$ are countable in $L_u$, and with the $<_L$-least surjections
$u$ would be countable in $L_{u^+}$. The second statement holds by the definition of $u^+$. (c) For infinite $x$,
$|F(x)| < x^+ < u$ by (L3) and (a); for finite $x$, $F(x)$ is a recursive order, so $|F(x)| < \omega^+ < u$.
\end{proof}

Two related results from the literature are not used in the proofs. Hoole (1986) studies an ordinal defined by
Slaman through recursion in a type-2 functional, and proves that it starts a gap in the constructible hierarchy
that lasts exactly up to the next admissible ordinal (Lemma 3 and Theorem 12 there). We expect that ordinal to be
$u$ but have not verified the identification.
Schlutzenberg (MathOverflow answer 381041, quoted by Schweber, Section 3.2) considers $\omega_1^{L_\kappa}$ for
$L_\kappa$ the least model of KP + ``$\omega_1$ exists'' and shows that it is the least ordinal at which a
pointwise definability property of the levels $L_{\alpha}$ fails. By Lemma 2.4 that model is $L_{u^+}$ and the
ordinal is $u$. His property concerns definability of the elements of $L_{\alpha}$; Theorem 3 concerns
$\Sigma^1_1$ classes of ordinals, and its threshold is $\mu_0 < u$.

Let $H$ be the set of elements of $L_{u^+}$ that are $\Sigma_1$-definable from $u$ in $L_{u^+}$, and let $\mu_0$ be
the least ordinal not in $H$. By (L1), $H$ is a $\Sigma_1$-elementary substructure of $L_{u^+}$ and there are a
unique ordinal $\gamma$ and a unique isomorphism $\pi : H \cong L_{\gamma}$.

\begin{statement}\head{Lemma 2.5} (a) $\mu_0 < u$, and $H \cap L_u = L_{\mu_0}$. So $\pi$ is the identity on
$L_{\mu_0}$ and $\pi(u) = \mu_0$.

(b) If $d < \mu_0$ is infinite then $d^+ < \mu_0$. So $\mu_0$ is a limit of admissible ordinals and a closure point
of the recursive dilators.
\end{statement}

\begin{proof} (a) $u$ is uncountable in $L_{u^+}$, so $\mu_0 < u$ by Lemma 2.3. Let $z \in H \cap L_u$. By Lemma
2.4 (a) $z$ is countable in $L_u$, and if $z$ is not empty the $<_L$-least surjection from $\omega$ onto $z$ is $\Sigma_1$-definable
from $z$, so it and its values lie in $H$. Hence $H \cap L_u$ is transitive, and its ordinals are those below
$\mu_0$. If $\rho < \mu_0$ then $L_{\rho} \in H$, so $L_{\rho} \subseteq H$. If $z \in H \cap L_u$, the least
$\rho$ with $z \in L_{\rho + 1}$ is in $H \cap u = \mu_0$. So $H \cap L_u = L_{\mu_0}$; here $\mu_0$ is a limit
because $H$ is closed under successor. The collapse is the identity on a transitive set and sends $u$ to the order
type of $H \cap u$. (b) $d^+ < u$ is $\Sigma_1$-definable from $d \in H$, so $d^+ \in H \cap u = \mu_0$; the rest
is as in Lemma 2.4 (c).
\end{proof}

\begin{statement}\head{Lemma 2.6} (a) Every element of $L_{\gamma}$ is $\Sigma_1$-definable from $\mu_0$ in
$L_{\gamma}$, and $\mu_0$ is uncountable in $L_{\gamma}$.

(b) $L_{\gamma + 1}$ contains a map from $\omega$ onto $L_{\gamma}$. So $\gamma$ is the least ordinal such that
$\mu_0$ is countable in $L_{\gamma + 1}$, and $\mu_0 < \gamma < \mu_0^+$.
\end{statement}

\begin{proof} (a) Definitions in $H$ are transported by $\pi$. That no surjection from $\omega$ onto $u$ exists is
a $\Pi_1$ fact of $L_{u^+}$, so it holds in $H$ and, transported, in $L_{\gamma}$. (b) $\gamma$ is a limit ordinal,
so $\Sigma_1$ satisfaction for $L_{\gamma}$ is definable over $L_{\gamma}$. Sending a $\Sigma_1$ formula to the
$<_L$-least solution of it with parameter $\mu_0$ in $L_\gamma$ is therefore a map in $L_{\gamma+1}$ from a set of formulas onto $L_{\gamma}$, by (a).
Since $\mu_0 < u$, $\mu_0$ is countable in $L_{\mu_0^+}$; as it is uncountable in $L_{\gamma}$, $\gamma < \mu_0^+$.
\end{proof}

\subsection{Definable well-orderings}

The ordinals $\delta_{\alpha}$ and $\dz{\alpha}$ are defined in Section 1.1. Both are at most $\alpha^+$, since a
relation definable over $L_{\alpha}$, as a set of ordered pairs, lies in $L_{\alpha+3}$ and (L3) applies.

\section{Trees with labels}

\subsection{The Kleene--Brouwer order of a labelled tree}

A \emph{pattern} of length $n$ consists of \emph{rows} $(k_i, m_i)$, $i < n$, with a \emph{tag} $k_i \in \omega$ and a
number $m_i \in \omega$ of \emph{slots}, together with a total preorder on the set of all slots
$\{(i, j) : i < n,\ j < m_i\}$. A \emph{pattern set} is a set $\mathcal S$ of patterns that contains the
empty pattern and is closed under restriction to the first $n'$ rows. For a linear order $b$, $\mathcal S(b)$ is the tree of all finite sequences
$((k_i, \vec a_i))_{i < n}$ with $\vec a_i \in b^{m_i}$ whose pattern, given by the tags, the lengths and the order
of the labels $\vec a_i$ in $b$, is in $\mathcal S$. Rows $(k, \vec a)$ are ordered by the code of $(k, m)$ and
then lexicographically. $P_{\mathcal S}(b)$ is the set of nodes of $\mathcal S(b)$ in the Kleene--Brouwer order:
$s$ is below $s'$ if $s$ properly extends $s'$, or $s$ has the smaller row at the first place where they differ.
The rank of a node $s$ is $\sup\{\operatorname{rk}(s') + 1 : s' \text{ a child of } s\}$, and $\infty$ if an infinite
branch passes through it; $\operatorname{rk} \mathcal S(b)$ is the rank of the root. Turning a tree into a functor of the order
from which its labels are taken is the standard way to obtain predilators (Girard 1981).

\begin{statement}\head{Lemma 3.1} If $\mathcal S$ is recursive then $P_{\mathcal S}$ is a recursive predilator, with
an index in $\Pre$ that is primitive recursive in an index of $\mathcal S$. For a well-order $b$, $P_{\mathcal S}(b)$
is a well-order if and only if $\mathcal S(b)$ has no infinite branch, and then
$|P_{\mathcal S}(b)| \ge \operatorname{rk} \mathcal S(b) + 1$.
\end{statement}

\begin{proof} The Kleene--Brouwer order is linear. A trace is a pattern in $\mathcal S$, its arity is the number of
classes of its preorder, and a node is a trace together with the increasing tuple of its distinct labels. Whether
one node is below another is determined by the two patterns and the order pattern of the two tuples, recursively.
An embedding $f$ acts by relabelling, which keeps patterns and the order of rows. If $f \le g$ pointwise and
$f \circ s \ne g \circ s$, the first row where they differ has the same tag in both and a smaller label in
$f \circ s$ at the first difference, so $f \circ s$ is below $g \circ s$. A branch gives a descending sequence of
its initial segments. If $b$ is a well-order and $\mathcal S(b)$ has no infinite branch, the rows are well-ordered
and the nodes extending a node $s$ form the ordered sum of the corresponding sets for the children of $s$, followed
by $s$; by induction on the rank this is a well-order of type at least $\operatorname{rk}(s) + 1$.
\end{proof}

\subsection{The tree of a $\Sigma^1_1$ sentence}

\textbf{Example 3.A.} Let $\theta$ be $\forall x\, \exists y\, (x < y)$, with a further binary relation symbol
$\tl$ about which $\theta$ says nothing. A node of the tree $\mathcal S(b)$ built below is a finite sequence of
rows. Row $i$ carries two labels $(x_i, w_i)$ from $b$ and a tag that says which pairs of labels of rows $\le i$
are in $\tl$. A fixed list schedules, for each $i$, one slot that occurs in a row $< i$ or is the slot of $x_i$,
and the node is admitted only if $w_i$ is greater than the label in the scheduled slot, the tags agree with each
other and on equal labels, and, for $i \ge 1$, the tag says $x_i \tl x_{i-1}$. Along an infinite branch every label is scheduled at some later row, which supplies a greater
label. So the set $A$ of labels of the branch has no greatest element, the relation read off from the tags is a
structure on $A$ satisfying $\theta$, and $x_0, x_1, \dots$ descends in it. At $b = 5$ there is no branch, since its
labels would form a nonempty finite set without a greatest element. At $b = \omega$ there is one: take $x_i = 0$
for all $i$, let $\tl$ hold between all pairs, and let each $w_i$ be the scheduled label plus one. If instead
$\theta$ also says that $\tl$ is the order $<$, then the tree for that sentence has no branch at $b = \omega$, but
every true node (in the sense of the proof below, for the structure $(\omega, <)$) whose last first label is $n$
has rank at least $n$, since it has a true child for each $m < n$; the root has such a child for every $n$, and so
the rank of the root is at least $\omega$, the height of $<$ on $\omega$. This rank bound is part (ii) of the
lemma.

\begin{statement}\head{Lemma 3.2} Let $\theta$ be a first-order sentence in a language $\{<, \vec R\}$ and $\tl$ a
binary relation symbol of $\vec R$. There is a recursive pattern set $\mathcal S = \mathcal S_{\theta, \tl}$, with an
index primitive recursive in $\theta$, such that for every linear order $b$:

(i) if $\mathcal S(b)$ has an infinite branch, then there are a nonempty $A \subseteq b$ and an expansion
$\mathfrak A$ of $(A, <)$ satisfying $\theta$ in which $\tl$ has an infinite descending sequence
$x_0 \mathrel{\triangleright} x_1 \mathrel{\triangleright} \dots$;

(ii) for every nonempty $A \subseteq b$ and every expansion $\mathfrak A$ of $(A, <)$ satisfying $\theta$,
$\operatorname{rk} \mathcal S(b) \ge \operatorname{ht}(\tl^{\mathfrak A})$, where
$\operatorname{ht}(\tl^{\mathfrak A}) = \sup\{\operatorname{rk}_{\tl}(x) + 1 : x \in A\}$ and
$\operatorname{ht} = \infty$ if $\tl^{\mathfrak A}$ is not well-founded.
\end{statement}

\begin{proof} By Skolem's normal form there is a sentence $\theta' = \forall \vec x\, \exists \vec y\, \theta_0$ in
a language with further relation symbols, $\theta_0$ quantifier-free, $|\vec x| = p$, $|\vec y| = q$, such that every
model of $\theta$ expands to a model of $\theta'$ and every model of $\theta'$ is a model of $\theta$. Fix a
primitive recursive list $(\mathbf v_i)_{i \in \omega}$ of $p$-tuples of slots in which every $p$-tuple occurs infinitely
often. Row $i$ of a pattern in $\mathcal S$ has $1 + q$ slots, with labels written $(x_i, \vec w_i)$, and its tag
codes a truth assignment $\Delta_i$ to all atomic formulas, in the relation symbols other than $<$, among the slots of
rows $\le i$. The pattern is in $\mathcal S$ if for every row $i$:
\begin{itemize}
\item[(a)] $\Delta_i$ extends $\Delta_{i-1}$, and gives equal values to atomic formulas that differ only by slots
  with equal labels;
\item[(b)] if every slot of $\mathbf v_i$ is the first slot of row $i$ or lies in a row $< i$, then
  $\theta_0(\mathbf v_i, \vec w_i)$ holds under the order pattern and $\Delta_i$;
\item[(c)] if $i \ge 1$, $\Delta_i$ makes $x_i \tl x_{i-1}$ true.
\end{itemize}
Condition (a) makes the assignments along a branch fit together into one structure on the set of labels. Condition
(b) tests one instance of $\forall \vec x\, \exists \vec y\, \theta_0$ per row, as soon as the arguments have
appeared, with the witnesses $\vec w_i$ supplied in the same row; since every tuple of slots is scheduled
infinitely often, every instance is tested on a branch. Condition (c) makes the first labels $x_0, x_1, \dots$ of
the rows a $\tl$-descending sequence; the first slot of a row has no other role, and it is this slot that ties the
rank of the tree to the height of $\tl$.

(i) Given a branch, let $A$ be the set of its labels and define the relations from the $\Delta_i$; by (a) this is
well defined. Every $p$-tuple from $A$ is the tuple of labels of some $\mathbf v_i$ with $i$ beyond the rows involved,
and (b) supplies witnesses. So the structure satisfies $\theta'$, hence $\theta$, and by (c) the $x_i$ descend.

(ii) Expand $\mathfrak A$ to $\mathfrak A' \models \theta'$. Call a node \emph{true} if its labels are in $A$ and
each $\Delta_i$ is the diagram of its slots in $\mathfrak A'$. Let $s$ be a true node, with last label $x$ in the
first slot of its last row if $s$ is not the root. For each $x' \in A$ with $x' \tl x$ (each $x' \in A$ if $s$ is the
root) $s$ has a true child whose last row has first label $x'$: its other labels are witnesses for $\mathbf v_i$ in
$\mathfrak A'$ if (b) asks for them, and arbitrary elements of $A$ otherwise, and its tag is the true diagram. By
induction along $\tl^{\mathfrak A}$, a true node with last first label $x$ has rank at least
$\operatorname{rk}_{\tl}(x)$, and the root has rank at least $\sup\{\operatorname{rk}_{\tl}(x) + 1\}$. If
$\tl^{\mathfrak A}$ has an infinite descending sequence, the true nodes along it form a branch.
\end{proof}

In (ii) the bound comes from the whole subtree of true nodes: the root has a true child for every $x \in A$, and
its rank is the supremum over all of them. One branch follows one descending sequence, which may be short.

When no relation
$\tl$ is given we add a symbol $\tl$ and the axiom $\forall x y\, (x \tl y)$ and write $\mathcal S_\theta$. Then
$\mathcal S_\theta(b)$ has an infinite branch if and only if some nonempty $A \subseteq b$ has an expansion
satisfying $\theta$. The suborders of an ordinal $b$ have exactly the order types $\le b$, so:
\begin{equation}
P_{\mathcal S_\theta}(b) \text{ is not well-ordered} \iff \text{some } \beta \text{ with } 0 < \beta \le b
\text{ lies in the class defined by } \theta . \tag{3.1}
\end{equation}

\subsection{Three codings}

The trees guess a structure on a subset $A$ of $b$ and test only quantifier-free conditions. The next lemma lists
what can be forced on a well-ordered $A$ by a first-order sentence.

\begin{statement}\head{Lemma 3.3} There are first-order sentences $\theta_c$, $\theta_L$ and, for each recursive
theory $T$ in the language $\{\in\}$, $\theta_T$, in finite relational languages extending $\{<\}$, with $\theta_L$
and $\theta_T$ containing $\theta_c$ as a conjunct and having a binary relation symbol $E$, such that for every
well-order $A$:

(C) $A$ has an expansion satisfying $\theta_c$ if and only if $A$ is infinite. In every such expansion, the elements
with finitely many predecessors form a copy $N$ of $(\omega, +, \cdot)$, and every finite sequence of elements of $A$ is
coded by an element of $A$. Hence for each $P \in \Pre$ there are formulas $\operatorname{dom}_P(x)$ and
$x <_P y$, primitive recursive in the index, and a map from $\{x : \operatorname{dom}_P(x)\}$ onto $P(A)$ such
that $x <_P y$ if and only if the image of $x$ is below the image of $y$.

(L) $A$ has an expansion satisfying $\theta_L$ if and only if $A$ is infinite, and in every such expansion $(A, E)$
is isomorphic to $(L_{\beta}, \in)$, where $\beta$ is the order type of $A$.

(K) In every expansion satisfying $\theta_T$, $(A, E) \models T$. Conversely, if $A$ is infinite, every structure
$(A, E)$ that is a model of $T$ expands to a model of $\theta_T$.

The sentences are primitive recursive in their data.
\end{statement}

The proof is in Appendix A. In (C), $n'$ is the
successor of $n \in N$.


In what follows, a single formula in the language $\{E\}$ may be conjoined to $\theta_T$ after relativisation, and
relation symbols of different sentences are kept disjoint.

\subsection{$\beta^-$-definable ordinals}

\begin{statement}\head{Lemma 3.4} (Ressayre, III.7, for predilators) (a) For $P \in \Pre$, the class of infinite
ordinals $\beta$ with $P(\beta)$ not well-ordered is $\Sigma^1_1$.

(b) For every $\Sigma^1_1$ class $X$ of infinite ordinals with least element $r$ there is $P \in \Pre$ such that,
for every ordinal $b$, $P(b)$ is well-ordered if and only if $b < r$.

So an infinite ordinal is $\beta^-$-definable if and only if it is the least infinite ordinal at which some
recursive predilator is not well-ordered.
\end{statement}

\begin{proof} (a) The class is defined by $\theta_c$ together with a unary function symbol $g$ and the axiom ``$g$ maps $N$ into
$\{x : \operatorname{dom}_P(x)\}$ and $g(n') <_P g(n)$ for all $n$''. (b) Let $\theta_X$ define $X$ and put
$P = P_{\mathcal S_\theta}$ for $\theta = \theta_X \wedge \theta_c$. By (3.1), $P(b)$ is not well-ordered if and
only if some infinite $\beta \le b$ is in $X$.
\end{proof}

\section{The criterion}

\textbf{Example 4.1.} Let $X = \{\omega\}$ and let $F$ be a recursive predilator that first fails at $\omega + 1$.
The tree $\mathcal S_\theta(b)$ for a sentence $\theta$ defining $X$ has an infinite branch at $b = \omega + 1$,
because the suborder $\omega$ of $b$ has type in $X$: trees are monotone in $b$, and no functorial tree has a
branch exactly when $b \in X$. So a tree that guesses a witness for ``$b \in X$'' together with a descending
sequence in $F(b)$ is not well-founded at $\omega + 1$, although $F$ is well-ordered on $X$. In the proof below the
descending sequence is taken in $F(A)$, for the subset $A \subseteq b$ on which the witness lives; this is achieved
by coding the elements of $F(A)$ as elements of $A$. A descending sequence in $F(b)$ would say nothing here,
because $F(b)$ can be ill-founded while $F(A)$, for the order $A$ that carries the witness, is a well-order.

\begin{proof}[Proof of Theorem 2 (a)] Left to right. Pick $D \in \Dil$ with $|F(\alpha)| \le |D(\alpha)|$. Let $X$
be the class defined by $\theta_c$ together with a unary function symbol $g$ and the axiom ``$g$ is a map from $\{x : \operatorname{dom}_F(x)\}$ to
$\{x : \operatorname{dom}_D(x)\}$ such that $x <_F y$ implies $g(x) <_D g(y)$''. Then $\alpha \in X$: let $g$ send a
code of $w \in F(\alpha)$ to a chosen code of the image of $w$ under a strictly increasing map $F(\alpha) \to D(\alpha)$;
that several codes may denote the same element does no harm. If $\beta \in X$, a descending
sequence in $F(\beta)$ would yield one in the well-order $D(\beta)$. So $F(\beta)$ is well-ordered.

Right to left. Let $\theta_X$ define $X$, and let $\theta$ be
\[ \theta_X \wedge \theta_c \wedge \forall x y\, (x \tl y \leftrightarrow x <_F y) . \]
Put $G = P_{\mathcal S_{\theta, \tl}}$, a recursive predilator by Lemma 3.1. If $G(b)$ is not well-ordered for
some ordinal $b$, Lemma 3.2 (i) gives $A \subseteq b$ with an expansion satisfying $\theta$ and an infinite
$<_F$-descending sequence. Then $A$ is a well-order, infinite by $\theta_c$, of a type $\beta \in X$, and by Lemma
3.3 (C) the sequence yields a descending sequence in $F(A) \cong F(\beta)$. This contradicts the hypothesis on $X$.
So $G \in \Dil$. Since $\alpha \in X$, $(\alpha, <)$ has an expansion satisfying $\theta$, and in it
$\operatorname{ht}(\tl) \ge |F(\alpha)|$, because every code has the rank of the element of $F(\alpha)$ that it denotes.
(Equality can fail: if $F(\alpha)$ is empty the height is $1$.) By Lemmas 3.2 (ii) and 3.1, $|G(\alpha)| \ge |F(\alpha)| + 1$. Hence
$|F(\alpha)| < |G(\alpha)| < \Xo(\alpha)$.
\end{proof}

The proof applies to every infinite $\alpha$, including $\alpha = \omega$ and uncountable $\alpha$.

\begin{proof}[Proof of Theorem 2 (b)] Suppose first that some $\beta^-$-definable ordinal is greater than $\alpha$,
and let $r = r(\alpha)$.

Right to left. Let $X$ be a $\Sigma^1_1$ class with $\alpha \in X \subseteq [0, r)$, and let $P \in \Pre$ with
$P(\alpha)$ well-ordered. If $P$ is well-ordered at every infinite ordinal it is well-ordered on $X$. Otherwise the
least infinite ordinal at which it is not well-ordered is $\beta^-$-definable by Lemma 3.4 (a) and greater than
$\alpha$, so it is at least $r$, and again $P$ is well-ordered on $X$. Theorem 2 (a) gives
$|P(\alpha)| < \Xo(\alpha)$. So $\Xp(\alpha) \le \Xo(\alpha)$.

Left to right. By Lemma 3.4 (b) there is $F \in \Pre$ that is well-ordered exactly below $r$. Then
$|F(\alpha)| < \Xp(\alpha) = \Xo(\alpha)$, and Theorem 2 (a) gives a $\Sigma^1_1$ class $X$ containing $\alpha$ on
which $F$ is well-ordered, so $X \subseteq [0, r)$.

If no $\beta^-$-definable ordinal is greater than $\alpha$, every $P \in \Pre$ with $P(\alpha)$ well-ordered is
well-ordered at every infinite ordinal by Lemma 3.4 (a), and Theorem 2 (a) applies with $X$ the class of all
infinite ordinals.
\end{proof}

\section{Terms over the next admissible set}

\begin{statement}\head{Lemma 5.1} Let $M = (A, E)$ be a model of KP + $V = L$ and $c \in A$ such that
$M \models$ ``$c$ is an ordinal'' and the $E$-predecessors of $c$ are well-ordered by $E$ in an infinite order type
$\beta$. Identify the well-founded part of $M$ with a transitive set. Then:

(a) $c = \beta$, $L_{\beta^+} \subseteq M$, and every element of $M$ that is in $M$ a member of an element of
$L_{\beta^+}$ belongs to $L_{\beta^+}$;

(b) every $\Pi_1$ formula with parameters in $L_{\beta^+}$ that holds in $M$ holds in $L_{\beta^+}$;

(c) if $\eta \le \beta$ is infinite and $t$ is a term with $t(\eta)$ defined, then $t^M(\eta) = t(\eta)$.
\end{statement}

\begin{proof} (a) $c$ is in the well-founded part and is the ordinal $\beta$. By (L4) the ordinals of the
well-founded part form an admissible ordinal greater than $\beta$, or $M$ is well-founded and isomorphic to an
admissible $L_\rho$ with $\rho > \beta$. In both cases, for $a < \beta^+$ the set that $M$ takes for $L_a$ is $L_a$,
by induction on $a$. (b) follows from (a). (c) The $<_L$-least pair $(w, v)$ of $L_{\eta^+}$ with
$\tau(w, v, \eta)$ lies in $L_{\beta^+}$ and satisfies $\tau$ in $M$, as $\tau$ is $\Delta_0$. The $<_L^M$-least
such pair of $M$ is not above it, so by (L1) and (a) it lies in the same level $L_a$, $a < \beta^+$, and it
satisfies $\tau$ in $L_{\beta^+}$. Hence it is $(w, v)$.
\end{proof}

Without ``$c$ is an ordinal'' the lemma fails: in $L_{\omega^+}$ the set $c = \omega \setminus \{0\}$ has predecessors
of order type $\omega$, but $c$ is not $\omega$.

Let $\theta_{\mathrm{KP}}$ be the sentence $\theta_T$ of Lemma 3.3 (K) for $T$ = KP + $V = L$, together with a
constant $c$, a function symbol $\iota$ and the statements that $(A, E) \models$ ``$c$ is an ordinal'' and that
$\iota$ is an isomorphism of $(A, <)$ onto the set of $E$-predecessors of $c$ ordered by $E$. For an infinite
ordinal $b$, the \emph{standard expansion} of $(b, <)$, used in Lemma 5.2, Corollary 5.3 and Lemma 7.1, is obtained from $L_{b^+}$ by a bijection with $b$, with
$c = b$ and true satisfaction; it satisfies $\theta_{\mathrm{KP}}$. Such a bijection exists also for uncountable $b$:
an elementary substructure of $L_{b^+}$ of the cardinality of $b$ that contains $b \cup \{b\}$ collapses to some
$L_{\rho}$ with $b < \rho \le b^+$ that satisfies KP, so $\rho = b^+$. In any expansion of a well-order $A$ of type
$\beta$ satisfying $\theta_{\mathrm{KP}}$, Lemma 5.1 applies to $(A, E)$ and $c$ with this $\beta$.

\begin{statement}\head{Lemma 5.2} Let $t$ be a term and $X$ a $\Sigma^1_1$ class of infinite ordinals. There is
$P = P_{t, X} \in \Pre$, with an index primitive recursive in $t$ and a sentence defining $X$, such that for every
ordinal $b$:

(i) if $t(\beta)$ is defined for every $\beta \le b$ with $\beta \in X$, then $P(b)$ is a well-order;

(ii) if $b \in X$, $t(b)$ is defined and is an ordinal, and $P(b)$ is a well-order, then $|P(b)| > t(b)$.
\end{statement}

\begin{proof} Let $\theta$ be the conjunction of a sentence $\theta_X$ defining $X$, $\theta_{\mathrm{KP}}$, the
statement that $(A, E) \models$ ``$t(c)$ is defined and is an ordinal'', and
\[ \forall x y\, \bigl(x \tl y \leftrightarrow x \E y \wedge (A, E) \models \text{``} y \le t(c) \text{''}\bigr), \]
where $t(c)$ is read as $t^M(c)$. Put $P = P_{\mathcal S_{\theta, \tl}}$. (i) An infinite branch at $b$ gives
$A \subseteq b$ of a type $\beta \le b$ with $\beta \in X$ and a model $M = (A, E)$ with an infinite $E$-descending
sequence of $M$-ordinals below $t^M(c)$. By Lemma 5.1 (a) and (c), $t^M(c) = t(\beta)$ is an ordinal of the
well-founded part of $M$, a contradiction. (ii) In the standard expansion of $(b, <)$ with a witness for $b \in X$,
$\operatorname{ht}(\tl) = t(b) + 1$. Apply Lemmas 3.2 (ii) and 3.1.
\end{proof}

\begin{statement}\head{Corollary 5.3} Let $\alpha$ be an infinite ordinal.

(a) $\Xp(\alpha) = \sup\{t(\alpha) : t$ a term with ordinal values that is defined at every infinite
$\beta \le \alpha\}$.

(b) $\Xo(\alpha) = \sup\{t(\alpha) : t$ a term with ordinal values that is defined at every element of some
$\Sigma^1_1$ class containing $\alpha\}$.

(c) If $\{\alpha\}$ is a $\Sigma^1_1$ class, then $\Xo(\alpha) = \Xp(\alpha)$ is the supremum of the ordinals that
are $\Sigma_1$-definable from $\alpha$ in $L_{\alpha^+}$ (Ressayre, IV.3).

(d) For every $\Pi_1$ formula $\varphi$ the class $\{\beta \ge \omega : L_{\beta^+} \models \varphi(\beta)\}$ is
$\Sigma^1_1$, defined by a sentence that is primitive recursive in $\varphi$ (Ressayre, III.5).
\end{statement}

\begin{proof} By Lemma 2.2 each value $|P(\alpha)|$ is $t_P(\alpha)$ for a term defined exactly where $P$ is
well-ordered; this gives $\le$ in (a) and, for $P \in \Dil$, in (b). For $\ge$ in (a) use $P_{t, X}$ with $X$ the
class of all infinite ordinals: by (i) it is well-ordered at $\alpha$, and by (ii) its value exceeds $t(\alpha)$.
For $\ge$ in (b), $P_{t, X}$ is well-ordered at every ordinal by (i), so it is a dilator, and (ii) applies at
$\alpha$. (c) is (b) with $X = \{\alpha\}$, together with $\Xp(\alpha) \le$ the supremum, by Lemma 2.2. (d) The
class is defined by $\theta_{\mathrm{KP}}$ together with $(A, E) \models \varphi(c)$: if $(\beta, <)$ has such an
expansion then $L_{\beta^+} \models \varphi(\beta)$ by Lemma 5.1 (b), and conversely the standard expansion
satisfies the sentence.
\end{proof}

Part (d) is the converse of (L5): together they say that, on countable infinite ordinals, the $\Sigma^1_1$ classes
are exactly the classes $\{\beta : L_{\beta^+} \models \varphi(\beta)\}$ with $\varphi$ a $\Pi_1$ formula.

By (a)--(c), $\Xo(\alpha)$, $\Xp(\alpha)$ and the supremum of the ordinals $\Sigma_1$-definable from $\alpha$ in
$L_{\alpha^+}$ are suprema of the values $t(\alpha)$ over the terms that are defined, respectively, on a
$\Sigma^1_1$ class containing $\alpha$, on $[\omega, \alpha]$, and at $\alpha$.

\section{Parameter-free definable well-orderings and Ressayre's functors}

\begin{statement}\head{Lemma 6.1} $\dz{\alpha} = \Xp(\alpha)$ for every infinite ordinal $\alpha$.
\end{statement}

\begin{proof} $\ge$. Let $P \in \Pre$ with $P(\alpha)$ a well-order. Recursive relations on $\omega$ are $\Sigma_1$
over $L_{\alpha}$, as computations are hereditarily finite. If $\alpha$ is a limit, the elements $(z, s)$ of
$P(\alpha)$ lie in $L_{\alpha}$, and the order of $P(\alpha)$ is $\Sigma_1$ over $L_{\alpha}$ without parameters. If
$\alpha = \lambda + n$ with $\lambda$ a limit and $0 < n < \omega$, represent an increasing tuple from $\alpha$ by a
finite set $s_0 \subseteq \lambda$ and a set $J \subseteq n$ standing for $\{\lambda + j : j \in J\}$. The relation
``$(z, J, s_0)$ and $(z', J', s_0')$ are such triples, with traces $z, z'$, $s_0 \cup s_0' \subseteq \nu$ for some limit ordinal $\nu$,
and the first is below the second in $P$, the points of $J$ and $J'$ being taken above $\nu$ in their natural
order'' is $\Sigma_1$ over $L_{\alpha}$ without parameters, because the limit ordinals in $L_{\lambda + n}$ are
those $\le \lambda$. It is a well-ordering of type $|P(\alpha)|$. So $\dz{\alpha} \ge |P(\alpha)|$.

$\le$. Let $R$ be a well-ordering defined over $L_{\alpha}$ by a $\Sigma_1$ formula $\varphi$ without parameters.
Let $\theta$ be $\theta_L \wedge \forall x y\, (x \tl y \leftrightarrow \varphi^E(x, y))$, where $\varphi^E$ is
$\varphi$ with $\in$ read as $E$, and $P = P_{\mathcal S_{\theta, \tl}}$. If $P(\alpha)$ were not a well-order,
Lemma 3.2 (i) and Lemma 3.3 (L) would give $\beta \le \alpha$ and an infinite sequence in $L_{\beta}$ that descends
in the relation defined by $\varphi$ over $L_{\beta}$. As $L_{\beta}$ is a transitive subset of $L_{\alpha}$ and
$\varphi$ is $\Sigma_1$, the sequence descends in $R$. So $P(\alpha)$ is a well-order. With $L_{\alpha}$ put on
$\alpha$ by a bijection, Lemmas 3.2 (ii) and 3.1 give $|P(\alpha)| \ge |R| + 1$.
\end{proof}

The proof of $\le$ uses the upward persistence of $\Sigma_1$ formulas. At limit $\alpha$ the relation used for
$\ge$ is $\Delta_1$, so there $\dz{\alpha}$ is also the supremum over the $\Delta_1$-definable well-orderings
without parameters; at successor $\alpha$ we have not determined this.

Ressayre's functors are different objects. Let $\mathcal R$ be the class of functors from the finite ordinals with
increasing maps to the finite ordinals that are primitive recursive on objects and maps, each extended to all
linear orders by direct limits (Ressayre, I.1--I.4). For an ordinal $\alpha$, $\PR(\alpha)$ is the supremum of
$|F(\alpha)|$ over the $F \in \mathcal R$ with $F(b)$ well-ordered for every ordinal $b$, and $\PRp(\alpha)$ is the
supremum of $|F(\alpha)|$ over the $F \in \mathcal R$ with $F(\alpha)$ well-ordered (Ressayre, I.10 and IV.0). A functor in
$\mathcal R$ need not preserve pullbacks: the functor with $F(0) = 0$ and $F(n) = 1$ for $n \ge 1$ does not
preserve the pullback of two disjoint one-element suborders. So the classes differ, and only the suprema are
compared. Appendix B proves the following (Lemma B.2): for every $F$ in $\mathcal R$ or in $\Pre$ there is a
$G$ in both such that for every infinite ordinal $b$, $G(b)$ is a well-order if and only if $F(b)$ is, and then
$|G(b)| > |F(b)|$.

\begin{proof}[Proof of Corollary 5] Let $\alpha$ be infinite. If $F \in \mathcal R$ is well-ordered at every
ordinal, the $G$ of Lemma B.2 is well-ordered at every infinite ordinal and finite at finite ones, so $G \in \Dil$
and $|F(\alpha)| < |G(\alpha)|$. If $D \in \Dil$, its $G$ is in $\mathcal R$, well-ordered at every ordinal, and
$|D(\alpha)| < |G(\alpha)|$. So $\PR (\alpha) = \Xo(\alpha)$. The same argument with ``well-ordered at
$\alpha$'' gives $\PRp (\alpha) = \Xp(\alpha)$, and Lemma 6.1 gives $\Xp(\alpha) = \dz{\alpha}$. The
rest restates Theorem 1 with Lemma 2.1; $\mu_0$ is Ressayre's $p_0$ by definition (proof of his Theorem IV.7).
\end{proof}

\section{$\beta$-definability and a self-referential predilator}

\begin{statement}\head{Lemma 7.1} Let $r$ be an infinite $\beta^-$-definable ordinal that is countable in
$L_{r^+}$. Then $r$ is $\beta$-definable.
\end{statement}

This is Ressayre's implication $(0)' \Rightarrow (0)$ of his Theorem IV.4, proved here for recursive predilators.

The proof finds an ordinal $\alpha_0 \le r$ with the same next admissible ordinal as $r$ that is the only element
of a $\Sigma^1_1$ class, and then defines $r$ from $\alpha_0$ inside $L_{r^+}$. The auxiliary ordinal $\alpha_1$ is
the left end of the interval of infinite ordinals with next admissible ordinal $r^+$: for infinite
$v < \alpha_1$ with $F(v)$ well-ordered, (L3) gives $|F(v)| < v^+ \le \alpha_1$, so only the interval $[\alpha_1, r]$ needs an
argument, and $\alpha_0$ is taken least in it so that Step 1 has no smaller competitor.

\begin{proof} Let $\alpha_1$ be the least infinite ordinal with $\alpha_1^+ = r^+$. If $v < \alpha_1$ is infinite
then $v^+ \le \alpha_1$: otherwise $v < \alpha_1 < v^+$ and $v^+ = \alpha_1^+ = r^+$. Let $\alpha_0$ be the least
$\beta^-$-definable ordinal in $[\alpha_1, r]$, and by Lemma 3.4 (b) let $F \in \Pre$ be well-ordered exactly below
$\alpha_0$. Note $\alpha_0^+ = r^+$.

\emph{Step 1: $|F(v)| < \alpha_0$ for every infinite $v < \alpha_0$.} Otherwise let $y$ be the least infinite
$v < \alpha_0$ with $|F(v)| \ge \alpha_0$. If $y < \alpha_1$ then $|F(y)| < y^+ \le \alpha_1 \le \alpha_0$ by (L3);
so $y \ge \alpha_1$. The composition $F \circ F$ is a recursive predilator and $F(F(v)) \cong F(|F(v)|)$. For
infinite $v < y$ this is a well-order, as $|F(v)| < \alpha_0$; at $y$ it is not, as $|F(y)| \ge \alpha_0$. So $y$
is the least infinite ordinal at which $F \circ F$ is not well-ordered, and by Lemma 3.4 it is a $\beta^-$-definable
ordinal in $[\alpha_1, \alpha_0)$. This contradicts the choice of $\alpha_0$.

\emph{Step 2: $\{\alpha_0\}$ is a $\Sigma^1_1$ class.} The class $Y$ consists of the infinite ordinals $\beta$ at
which $F$ is not well-ordered and such that $F(v)$ embeds in $\beta$ for every infinite $v < \beta$. Formally, $Y$ is
defined by $\theta_c$ together with a unary function symbol $g$, a binary function symbol $h$ and the axioms
``$g$ maps $N$ into $\{x : \operatorname{dom}_F(x)\}$ with $g(n') <_F g(n)$ for all $n$'', and ``for every $v$ not
in $N$, $h(v, \cdot)$ maps $\{x : \operatorname{dom}_F(x),\ x_1, x_2, \dots < v\}$ into $A$ and $x <_F y$ implies
$h(v, x) < h(v, y)$''. Let $\beta \in Y$. Then $\beta$ is infinite and $F(\beta)$ is not well-ordered, so
$\beta \ge \alpha_0$. If $\beta > \alpha_0$, the clause for $v = \alpha_0$ maps $F(\alpha_0)$ increasingly into the
well-order $\beta$, which is impossible. So $\beta = \alpha_0$. And $\alpha_0 \in Y$ by Step 1.

\emph{Step 3: $r = t(\alpha_0)$ for a term $t$.} Since $\alpha_0 \le r$ and $r$ is countable in
$L_{r^+} = L_{\alpha_0^+}$, $\alpha_0$ is countable in $L_{\alpha_0^+}$. By Lemma 2.3 every element of
$L_{\alpha_0^+}$, in particular $r$, is $\Sigma_1$-definable from $\alpha_0$ in $L_{\alpha_0^+}$.

\emph{Step 4: $\{r\}$ is a $\Sigma^1_1$ class.} Consider the sentence consisting of $\theta_{\mathrm{KP}}$; a
constant $c_0$ with $(A, E) \models$ ``$c_0$ is an ordinal, $c_0 \le c$, and $t(c_0) = c$''; and a sentence
defining $Y$, in new symbols, relativised to the initial segment $\{a : \iota(a) \E c_0\}$ of $A$. Let a well-order
$A$ of type $\beta$ have such an expansion $M$. By Lemma 5.1 (a), $c = \beta$ and $c_0$ is an ordinal
$\eta \le \beta$; the initial segment has type $\eta$, so $\eta \in Y$ and $\eta = \alpha_0$. By Lemma 5.1 (c),
$t^M(\alpha_0) = t(\alpha_0) = r$. So $\beta = r$. Conversely the standard expansion of $(r, <)$ with
$c_0 = \alpha_0$, and a witness for $\alpha_0 \in Y$ on the initial segment $\alpha_0$ of $r$, satisfies the
sentence; here $t^M(\alpha_0) = r$ because $L_{r^+} = L_{\alpha_0^+}$.
\end{proof}

\begin{statement}\head{Lemma 7.2} (Ressayre, Remark IV.5) $u$ is $\beta$-definable.
\end{statement}

\begin{proof} Let $\varphi(v)$ be the conjunction of: (a) $v$ is an ordinal and there is no surjection from
$\omega$ onto $v$; (b) $L_v$ is admissible; (c) for every infinite ordinal $v' < v$ and every $z$, if $z = L_\rho$
for the least admissible $\rho > v'$, then $z$ contains a surjection from $\omega$ onto $v'$. By (L1), and because
satisfaction in a set is $\Delta_1$, $\varphi$ is equivalent to a $\Pi_1$ formula in every admissible set. Let
$\beta$ be infinite. If $\beta < u$ then (a) fails in $L_{\beta^+}$. If $\beta = u$ then (a)--(c) hold in
$L_{u^+}$ by Lemma 2.4. If $u < \beta < u^+$ then (b) fails by Lemma 2.4 (b). If $\beta \ge u^+$ then (c) fails in
$L_{\beta^+}$ for $v' = u$ and $z = L_{u^+}$. So $\{\beta \ge \omega : L_{\beta^+} \models \varphi(\beta)\} = \{u\}$,
and Corollary 5.3 (d) applies.
\end{proof}

\begin{statement}\head{Lemma 7.3} Let $r$ be an infinite $\beta$-definable ordinal, and let $Z \subseteq \omega$ be
$\Sigma_1$-definable from $r$ in $L_{r^+}$ with $\Dil \subseteq Z$. Then there is $e \in Z \cap \Pre$ such that, for
every ordinal $b$, $P_e(b)$ is a well-order if and only if $b < r$.
\end{statement}

\begin{proof} Let $\theta_r$ define $\{r\}$, and let $\chi(z, v)$ be a $\Sigma_1$ formula such that $Z$ is the
unique $z$ with $L_{r^+} \models \chi(z, r)$. For a $\Sigma_1$ formula $\psi(v)$ let
\[ Y_\psi = \{\beta \ge \omega : \beta = r \text{ and } L_{\beta^+} \models \neg\psi(\beta)\} . \]
By Corollary 5.3 (d), $Y_\psi$ is defined by a sentence $\theta_\psi$ that contains $\theta_r \wedge \theta_c$ and is primitive recursive in $\psi$. Let
$p(\psi) \in \Pre$ be the index of $P_{\mathcal S_{\theta_\psi}}$ given by Lemma 3.1; $p$ is primitive recursive.
By (3.1):
\begin{itemize}
\item if $L_{r^+} \models \psi(r)$, then $Y_\psi = \emptyset$, $P_{p(\psi)}$ is well-ordered at every ordinal, and
  $p(\psi) \in \Dil$;
\item if $L_{r^+} \models \neg\psi(r)$, then $Y_\psi = \{r\}$ and $P_{p(\psi)}(b)$ is a well-order exactly for
  $b < r$.
\end{itemize}
The graph of $p$ is $\Sigma_1$ over $L_{\omega}$, so by the diagonal lemma there is a $\Sigma_1$ formula $\psi$ such
that in every admissible set
\[ \psi(v) \iff \exists z\, \bigl(\chi(z, v) \wedge p(\ulcorner\psi\urcorner) \notin z\bigr) . \]
If $L_{r^+} \models \psi(r)$, then $p(\psi) \notin Z$ and $p(\psi) \in \Dil$, contradicting $\Dil \subseteq Z$. So
$L_{r^+} \models \neg\psi(r)$. Since $Z$ satisfies $\chi$, $p(\psi) \in Z$, and $e = p(\psi)$ is as required.
\end{proof}

The lemma uses that $\Dil$ contains every index of a recursive dilator: the hypothesis is applied to the particular
index $p(\psi)$.

\section{Agreement below $\mu_0$ and the values at $\mu_0$}

For an infinite ordinal $d$ let
\[ Z_d = \{e \in \Pre : P_e(d) \text{ is a well-order of type} < \Xo(d)\} . \]

\begin{statement}\head{Lemma 8.1} $\Dil \subseteq Z_d \in L_{d^+ + 1}$, and if $\Xo(d) < d^+$ then
$Z_d \in L_{d^+}$.
\end{statement}

\begin{proof} $\Dil \subseteq Z_d$ because the supremum $\Xo(d)$ is not attained. If $\Xo(d) = d^+$, then by (L3)
$Z_d$ is the set of $e \in \Pre$ such that $L_{d^+}$ contains an isomorphism of $P_e(d)$ onto an ordinal, which is
definable over $L_{d^+}$. Let $\zeta = \Xo(d) < d^+$. Define by $\Sigma$-recursion in $L_{d^+}$, for $e \in \Pre$
and $\xi \le \zeta$: $\nx(e, \xi)$ is the least element of $P_e(d)$ that is not among the $\nx(e, \xi')$, $\xi' < \xi$,
if there is a least one, and a fixed marker otherwise. The restriction of $\nx$ to $\Pre \times (\zeta + 1)$ is an
element of $L_{d^+}$, and $e \in Z_d$ if and only if for some $\xi < \zeta$ no marker occurs below $\xi$ and the
values $\nx(e, \xi')$, $\xi' < \xi$, exhaust $P_e(d)$. So $Z_d \in L_{d^+}$ by $\Delta_0$-separation.
\end{proof}

\begin{statement}\head{Lemma 8.2} Every infinite ordinal $d < \mu_0$ belongs to a $\Sigma^1_1$ class all of whose
elements are below $u$.
\end{statement}

\begin{proof} By Lemmas 8.1 and 2.5, $Z_d \in L_{d^+ + 1} \subseteq L_{\mu_0} \subseteq H$, so $Z_d$ is
$\Sigma_1$-definable from $u$ in $L_{u^+}$. By Lemmas 7.2 and 7.3 there is $e \in Z_d$ such that $P_e$ is
well-ordered exactly below $u$. Since $|P_e(d)| < \Xo(d)$, Theorem 2 (a) gives a $\Sigma^1_1$ class containing $d$
on which $P_e$ is well-ordered.
\end{proof}

\begin{proof}[Proof of Theorem 1 (i)] Let $d < \mu_0$ be infinite. By Lemmas 7.2 and 2.5, $u$ is a
$\beta^-$-definable ordinal above $d$, so $r = r(d)$ exists and $r \le u$. If $r = u$, Lemma 8.2 and Theorem 2 (b)
give $\Xp(d) = \Xo(d)$. Let $r < u$. If $\Xo(d) = d^+$ then $\Xp(d) = \Xo(d)$ by (L3). Otherwise
$Z_d \in L_{d^+} \subseteq L_{r^+}$ by Lemma 8.1. By Lemma 2.4 (a), $r$ is countable in $L_{r^+}$. So $r$ is
$\beta$-definable by Lemma 7.1, and $Z_d$ is $\Sigma_1$-definable from $r$ in $L_{r^+}$ by Lemma 2.3. Lemma 7.3
gives $e \in Z_d$ such that $P_e$ is well-ordered exactly below $r$, Theorem 2 (a) gives a $\Sigma^1_1$ class
containing $d$ and contained in $[0, r)$, and Theorem 2 (b) gives $\Xp(d) = \Xo(d)$.
\end{proof}

\begin{statement}\head{Lemma 8.3} Every $\Sigma^1_1$ class that contains $\mu_0$ contains $u$. Consequently, if
$F_0 \in \Pre$ first fails at $u$, then $|F_0(\mu_0)| \ge \Xo(\mu_0)$, and $\Xp(\mu_0) > \Xo(\mu_0)$.
\end{statement}

\begin{proof} Let $X$ be a $\Sigma^1_1$ class and $\varphi$ a $\Pi_1$ formula for it as in (L5); $\mu_0$ and $u$
are countable. Suppose $u \notin X$. Then $L_{u^+} \models \exists w\, \chi(w, u)$ for a $\Delta_0$ formula $\chi$
equivalent to the matrix of $\neg\varphi$. There is such a $w$ in $H$, and $L_{\gamma} \models \chi(\pi(w), \mu_0)$.
Since $L_{\gamma}$ is transitive and $\gamma < \mu_0^+$ (Lemma 2.6), $L_{\mu_0^+} \models \neg\varphi(\mu_0)$, so
$\mu_0 \notin X$. For the second statement, $F_0(\mu_0)$ is well-ordered because $\mu_0 < u$. If
$|F_0(\mu_0)| < \Xo(\mu_0)$, Theorem 2 (a) gives a $\Sigma^1_1$ class containing $\mu_0$ on which $F_0$ is
well-ordered; it contains $u$, where $F_0$ is not well-ordered. Finally $\Xp(\mu_0) > |F_0(\mu_0)|$. Such an $F_0$
exists by Lemmas 7.2 and 3.4 (b).
\end{proof}

\begin{statement}\head{Lemma 8.4} $\Xo(\mu_0) = \gamma$.
\end{statement}

\begin{proof} $\le$. Let $D \in \Dil$. The order $D(u)$, its order type $\xi_D < u^+$ and the isomorphism
$f_D : D(u) \cong \xi_D$ are $\Sigma_1$-definable from $u$ in $L_{u^+}$ by (L3), so they lie in $H$. The statement
``$f$ is an isomorphism of $D(v)$ onto the ordinal $\xi$'' is $\Delta_0$ in a transitive set containing the
computation of $D(v)$, so $\pi(f_D)$ is an isomorphism of $D(\mu_0)$ onto $\pi(\xi_D) < \gamma$. Hence
$|D(\mu_0)| = \pi(\xi_D) < \gamma$.

$\ge$. Let $\eta < \gamma$, and let $\eta^* \in H$ be the ordinal with $\pi(\eta^*) = \eta$. Then
$\eta^* = t(u)$ for a term $t$. By Lemma 7.2, $\{u\}$ is a $\Sigma^1_1$ class, and $t$ is defined on it. By
Corollary 5.3 (b) there is $D \in \Dil$ with $|D(u)| > \eta^*$, and then $|D(\mu_0)| = \pi(\xi_D) > \eta$, as $\pi$
preserves the order of the ordinals in $H$.
\end{proof}

\begin{statement}\head{Lemma 8.5} Let $d$ be an infinite ordinal with $\Xo(d) < d^+$, and let $x \ge d$ be such that
every $\Sigma^1_1$ class containing $d$ has an element $\ge x$. Then $\Xo(x) < x^+$.
\end{statement}

\begin{proof} By Lemma 8.1, $Z_d \in L_{d^+} \subseteq L_{x^+}$. Let $e \in Z_d$. Since $|P_e(d)| < \Xo(d)$, by
Theorem 2 (a) there is a $\Sigma^1_1$ class containing $d$ on which $P_e$ is well-ordered.
It has an element $x' \ge x$, and $P_e(x)$ is a suborder of $P_e(x')$. So $P_e(x)$ is a well-order for every
$e \in Z_d$, and by (L3) its order type is in $L_{x^+}$. The map $e \mapsto |P_e(x)|$ on $Z_d$ is $\Sigma_1$ over
$L_{x^+}$, so by $\Sigma$-collection its values are bounded below $x^+$. Since $\Dil \subseteq Z_d$, the bound
applies to $\Xo(x)$.
\end{proof}

\begin{statement}\head{Lemma 8.6} No $\beta^-$-definable ordinal lies in $[\mu_0, u)$. Hence $r(x) = u$ for every
$x \in [\mu_0, u)$, the recursive predilators that are well-ordered at $x$ are those that are well-ordered at every
ordinal below $u$, and $\Xp$ is nondecreasing on $[\mu_0, u)$.
\end{statement}

\begin{proof} Let $r < u$ be infinite and $\beta^-$-definable, and by Lemma 3.4 let $P \in \Pre$ first fail at $r$.
For $\beta < u$, $P(\beta)$ is a well-order if and only if $L_{\beta^+}$ contains an isomorphism of it onto an
ordinal, and $L_{\beta^+} \in L_u$. So $r$ is the least infinite $\beta$ such that $L_{\beta^+}$ contains no such
isomorphism, which is a $\Sigma_1$ definition of $r$ in $L_{u^+}$. Hence $r \in H \cap u = \mu_0$. The rest follows
from Lemma 7.2 and Section 2.1.
\end{proof}

\begin{proof}[Proof of Theorem 1 (ii) and (iii)] $\mu_0$ is a closure point by Lemma 2.5, and $\Xo(\mu_0) = \gamma$
by Lemma 8.4; so $\Xo(x) = \gamma$ for $\mu_0 \le x < \gamma$ by Lemma 2.1. Recall $\mu_0 < \gamma < \mu_0^+ < u$.

\emph{$\Xp(\mu_0) \ge \Xp(\gamma)$.} Let $F_0$ be as in Lemma 8.3, so $|F_0(\mu_0)| \ge \gamma$. Let $P \in \Pre$ be
well-ordered at $\gamma$. By Lemma 8.6, $P$ is well-ordered below $u$. The composition $P \circ F_0$ is a recursive
predilator, and for $x < u$ it is well-ordered at $x$, since $|F_0(x)| < u$ by Lemma 2.4 (c). So
$\Xp(\mu_0) > |P(F_0(\mu_0))|$ ${}= |P(|F_0(\mu_0)|)| \ge |P(\gamma)|$.

\emph{$\Xp(\gamma) = \Xo(\gamma)$.} We need a $\Sigma^1_1$ class that contains $\gamma$ and no ordinal $\ge u$. The
class $X$ below says that every infinite $\alpha' \le \alpha$ is countable early: by level $\alpha + 3$ at the
latest, and already in its own next admissible set whenever that set lies below $\alpha$. The shift by $3$ leaves
room for the surjections onto the ordinals in $[\mu_0, \gamma]$, which first appear just above $L_{\gamma}$. The
clause about $\alpha'^+$ excludes the ordinals $\alpha \ge u^+$, for which $L_{\alpha+3}$ does contain a surjection
onto $u$.

Let $X$ be the class of infinite $\alpha$ such that for every infinite
$\alpha' \le \alpha$ there is a surjection $f$ from $\omega$ onto $\alpha'$ in $L_{\alpha + 3}$, and
$f \in L_{\alpha'^+}$ if $\alpha'^+ \le \alpha$. This class is $\Sigma^1_1$: for an infinite well-order $(A, <)$ of
type $\alpha$, moving its first three elements to the top gives a definable well-order of type $\alpha + 3$, to which
Lemma 3.3 (L) applies. The condition is first-order in the resulting expansion: the level predicate $B_a$ and the
satisfaction predicate of Appendix A are correct for every positive index used here, the indices are matched with the ordinals
of the copy by an order isomorphism, ``$L_{\rho}$ is admissible'' is tested by checking the recursive list of axioms
of KP in the satisfaction predicate for $B_\rho$, and the second clause requires $B_\rho(f)$ for the least admissible
$\rho$ with $\alpha' < \rho \le \alpha$. No $\alpha \ge u$
is in $X$: take $\alpha' = u$; if $u^+ \le \alpha$ there is no surjection in $L_{u^+}$, and otherwise
$L_{\alpha+3} \subseteq L_{u^+}$. And $\gamma \in X$: for $\alpha' < \mu_0$ use Lemma 2.4 (a) and
$\alpha'^+ < \mu_0$; for $\mu_0 \le \alpha' \le \gamma$, Lemma 2.6 (b) gives a surjection in $L_{\gamma+2}$, and
$\alpha'^+ = \mu_0^+ > \gamma$. So $\gamma \in X \subseteq [0, u)$ and $r(\gamma) = u$, and Theorem 2 (b) applies.

\emph{$\Xo(\gamma) < \mu_0^+$.} Apply Lemma 8.5 with $d = \mu_0$ and $x = \gamma$: $\Xo(\mu_0) = \gamma < \mu_0^+$,
and every $\Sigma^1_1$ class containing $\mu_0$ contains $u > \gamma$ by Lemma 8.3. So
$\Xo(\gamma) < \gamma^+ = \mu_0^+$.

For $\mu_0 \le x < \gamma$, Lemma 8.6 now gives
$\Xo(\gamma) = \Xp(\gamma) \le \Xp(\mu_0)$ ${}\le \Xp(x) \le \Xp(\gamma)$, and $\gamma < \Xo(\gamma)$.
\end{proof}

\section{Characterisations of $\mu_0$}

\begin{proof}[Proof of Theorem 3] (a) is Lemma 8.3. (b) By (a), $\mu_0$ has the property. A finite $d$ is the only
element of a $\Sigma^1_1$ class, and an infinite $d < \mu_0$ does not have the property by Lemma 8.2. (c) By (a),
$\mu_0$ has the property with $y = u$. Let $d < y$ be such that $y$ belongs to every $\Sigma^1_1$ class containing
$d$; then $d$ is infinite. We show $d \ge \mu_0$.

By Corollary 5.3 (d), every $\Pi_1$ formula true of $d$ in $L_{d^+}$ is true of $y$ in $L_{y^+}$. So every
$\Sigma_1$ formula true of $y$ in $L_{y^+}$ is true of $d$ in $L_{d^+}$. Let $H_y$ be the set of elements of
$L_{y^+}$ that are $\Sigma_1$-definable from $y$, and $\pi_y : H_y \cong L_{\gamma_y}$ its collapse. The elements
of $H_y$ are the values $t(y)$ of terms. Put $j(t(y)) = t(d)$. Each of the statements ``$t(y)$ is
defined'', $t(y) = t'(y)$, $t(y) \ne t'(y)$, $t(y) \in t'(y)$, $t(y) \notin t'(y)$ and ``$t(y)$ is an ordinal'',
when true, is a $\Sigma_1$ fact about $y$ in $L_{y^+}$: it says that there are least pairs for the terms
involved whose second components are so related. Hence it holds of $d$. So $t(d)$ is defined whenever $t(y)$ is;
$j$ is well defined by the case of $=$ and injective by the case of $\ne$; it preserves $\in$ and $\notin$; and it
maps ordinals to ordinals. So $j$ is strictly increasing on the ordinals of $H_y$, and
$j \circ \pi_y^{-1}$ is a strictly increasing map from $\gamma_y$ into the ordinals, and
$d = j(y) \ge \pi_y(y)$.

If $y < u$, then $y$ is countable in $L_{y^+}$, so $y \subseteq H_y$ by Lemma 2.3 and $\pi_y(y) = y > d$, a
contradiction. So $y \ge u$. Then $u \in H_y$: if $y < u^+$, $u$ is the greatest admissible ordinal $\le y$ by Lemma
2.4 (b); if $y \ge u^+$, $u$ is the least ordinal $\alpha$ such that $L_{\alpha^+}$ has no surjection from $\omega$
onto $\alpha$, a $\Sigma_1$ definition in $L_{y^+}$. In both cases $L_{u^+}$ is $L_{y^+}$ or an element of $H_y$, so
every ordinal below $\mu_0$, being $\Sigma_1$-definable from $u$ in $L_{u^+}$, is in $H_y$. Now $\pi_y(u)$ is the order type of the set of ordinals of $H_y$ below $u$, which contains every
ordinal below $\mu_0$, and $u \le y$. Hence $\pi_y(y) \ge \pi_y(u) \ge \mu_0$.
\end{proof}

\section{The next closure point: four candidates}

\begin{proof}[Proof of Theorem 4] (a) Let $\alpha$ be infinite. By the first hypothesis $f_{\mathcal O}(\alpha)$ is a
closure point of the recursive dilators above $\alpha$: for finite $\beta$, $|D(\beta)| \le |D(\omega)|$. So
$f_{\mathcal O}(\alpha) \ge \Xo(\alpha)$. Let $\omega \le \beta < \Xo(\alpha)$ and $P \in \mathcal O_\beta$. By the
second hypothesis and Theorem 2 (a), $|P(\beta)| < \Xo(\beta) \le \Xo(\alpha)$, using Lemma 2.1. So $\Xo(\alpha)$
has the defining property of $f_{\mathcal O}(\alpha)$.

(b) Every recursive dilator is in every $\mathcal O^*_\beta$. By the definition of $\Xp$,
$f_{\mathcal O^*}(\alpha)$ is the least $\kappa > \alpha$ with $\Xp(\beta) \le \kappa$ for all infinite
$\beta < \kappa$. For infinite $\alpha < \mu_0$, $\Xo(\alpha) \le \mu_0$ since $\mu_0$ is a closure point, so for
infinite $\beta < \Xo(\alpha)$ Theorem 1 (i) gives $\Xp(\beta) = \Xo(\beta) \le \Xo(\alpha)$, and
$f_{\mathcal O^*}(\alpha) = \Xo(\alpha)$. And $f_{\mathcal O^*}(\mu_0) \ge \Xp(\mu_0) > \Xo(\mu_0)$ by Theorem 1 (ii).
Finally, every $P \in \Pre$ well-ordered at an infinite $\beta$ is well-ordered on the class of infinite ordinals at
which $P$ is well-ordered, which contains $\beta$ and is $\Pi^1_1$ by Lemma 3.4 (a). So the $\Pi^1_1$ form of the
second hypothesis holds for $\mathcal O^*$.
\end{proof}

\begin{proof}[Proof of Proposition 6] We use the following cited results. (1) $\Xo(\alpha) = \alpha^+$ for every
infinite $\alpha < \sigma$: every ordinal $\le \alpha$ is $\beta$-definable by the definition of $\sigma$, so
Girard's theorem in the form of Ressayre's Theorem I.6 applies; Corollary 5 translates it to $\Xo$. (2)
$\Xp(\sigma) = \Xo(\sigma) < \sigma^+$ (Ressayre, the example after IV.6), and $\sigma$ is a limit of admissible
ordinals (Ressayre, proof of III.10). (3) For a real $y$, the supremum $\omega_1^y$ of the order types of the
well-orderings recursive in $y$ is an admissible ordinal (the relativised form of the admissibility of
the least non-recursive ordinal; Sacks).

\emph{$\sigma < \mu_0$.} By (1) and Theorem 1 (ii), $\sigma \le \mu_0$, as $\Xo(\mu_0) = \gamma < \mu_0^+$. By (2)
and Theorem 1 (ii), $\sigma \ne \mu_0$. So $\sigma + \omega < \mu_0$ by Lemma 2.5 (b).

\emph{$\dz{\sigma + \omega} = \Xo(\sigma) < \sigma^+$.} The recursive dilator $b \mapsto b + \omega + 1$ shows
$\sigma + \omega < \Xo(\sigma)$, so $\Xo(\sigma + \omega) = \Xo(\sigma)$ by Lemma 2.1. Now use Corollary 5, Theorem 1 (i) and (2).

\emph{$\delta_{\sigma + \omega} = \sigma^+$.} Every infinite $\alpha < \sigma$ is the only element of a $\Sigma^1_1$
class. By (L5), $\alpha$ is then the only countable infinite $\beta$ with $L_{\beta^+} \models \varphi(\beta)$, for
some $\Pi_1$ formula $\varphi$, and $L_{\alpha^+} \in L_{\sigma}$ by (2). Satisfaction in $L_{\beta^+}$ for
$\beta < \sigma$ is computed in a larger admissible level below $\sigma$, so the map that sends $\varphi$ to the
least infinite $\beta < \sigma$ with $L_{\beta^+} \models \varphi(\beta)$, where there is one, is definable over
$L_{\sigma}$; its range is $[\omega, \sigma)$. Together with separate codes for the finite ordinals this gives a map
in $L_{\sigma + 1}$ from a set of natural numbers onto $\sigma$. Comparing two codes by their values gives a
prewellordering $y$ of a set of natural numbers, of type $\sigma$, with $y \in L_{\sigma + 2}$. Its restriction to
the least code in each class is a well-ordering of type $\sigma$ that is recursive in $y$. Every well-ordering recursive in
$y$ is $\Delta_1$-definable over $L_{\sigma + \omega}$ with the parameter $y$, so
$\delta_{\sigma + \omega} \ge \omega_1^y$. By (3), $\omega_1^y$ is an admissible ordinal greater than $\sigma$, so
$\omega_1^y \ge \sigma^+$. And $\delta_{\sigma + \omega} \le (\sigma + \omega)^+ = \sigma^+$.

Finally $\sigma + \omega < \delta_{\sigma}$: the order of the set $(\{0\} \times \sigma) \cup (\{1\} \times \omega) \cup \{(2, 0)\}$
by first and then second coordinates has type $\sigma + \omega + 1$ and is $\Delta_1$-definable over $L_{\sigma}$. 
\emph{$\delta_{\sigma} < \sigma^+$.} We eliminate parameters. Each infinite $\beta < \sigma$ is the only element of
a $\Sigma^1_1$ class, so by (L5) it is the only countable infinite ordinal with
$L_{\beta^+} \models \varphi(\beta)$ for some $\Pi_1$ formula $\varphi$. Hence $\beta$ is $\Sigma_1$-definable
over $L_{\sigma}$ without parameters, as the ordinal for which some set is the next admissible level $L_{\beta^+}$
and satisfies $\varphi(\beta)$; that this set is $L_{\beta^+}$ is certified in a larger admissible level below
$\sigma$ (compare Ressayre, III.6 (a)). Finite ordinals are definable outright. Now let $p \in L_{\sigma}$, and
choose an infinite $\eta < \sigma$ with $p \in L_{\eta^+}$. By (L2), $p$ is $\Sigma_1$-definable in $L_{\eta^+}$
from ordinals $\le \eta$; substituting the definitions of these ordinals and of $\eta$, and certifying
$L_{\eta^+}$ in the same way, shows that $p$ is $\Sigma_1$-definable over $L_{\sigma}$ without parameters. So a
well-ordering that is $\Delta_1$-definable over $L_{\sigma}$ with the parameter $p$ is $\Sigma_1$-definable over
$L_{\sigma}$ without parameters: quantify existentially over the set defined by the definition of $p$. Hence
$\delta_{\sigma} \le \dz{\sigma} = \Xp(\sigma) < \sigma^+$ by Corollary 5 and (2).
\end{proof}

By (1) and (2) in this proof, $\alpha^+$ is the least closure point above $\alpha$ for every infinite
$\alpha < \sigma$ and not for $\alpha = \sigma$; this is the first row of the table in Section 1.2. We do not
determine the least infinite $\alpha$ with $\delta_{\alpha} \ne \Xo(\alpha)$.

\section{Remarks and questions}

\begin{enumerate}
\item \textbf{An embedding.} Let $\hat u$ be the supremum of the ordinals in $H$. Then $\Xo(u) = \hat u$ by Lemma
  7.2 and Corollary 5.3 (c), and $H$ is a $\Sigma_1$-elementary substructure of $L_{\hat u}$, because
  $H \subseteq L_{\hat u} \subseteq L_{u^+}$. So $\pi^{-1} : L_{\Xo(\mu_0)} \to L_{\Xo(u)}$ is a cofinal
  $\Sigma_1$-elementary embedding with critical point $\mu_0$ that sends $\mu_0$ to $u$. Moreover $L_{\mu_0}$ is
  an elementary substructure of $L_u$: if $L_u \models \exists z\, \varphi(z, q)$ with $q \in L_{\mu_0}$, the
  $<_L$-least such $z$ in $L_u$ is $\Sigma_1$-definable in $L_{u^+}$ from $q$ and $L_u$, which are in $H$, so it
  lies in $H \cap L_u = L_{\mu_0}$. These are standard facts about $\Sigma_1$ hulls; only the identification of the
  two heights with values of $\Xo$ uses this paper.
\item \textbf{$\mu_0$ is not the least ordinal with $L_{\alpha} \prec L_u$.} That ordinal is $\Sigma_1$-definable
  from $u$ in $L_{u^+}$, hence below $\mu_0$.
\item \textbf{Nearby statements that fail.} $\Xp = \Xo$ does not hold at every closure point of the recursive
  dilators ($\mu_0$), nor at every $\alpha$ for which $r(\alpha)$ is $\beta$-definable ($\mu_0$ again, by Lemma
  7.2). It does not fail only at $\mu_0$ below $u$: it fails on $[\mu_0, \gamma)$, and by Ressayre's Theorem IV.8
  (c) it fails cofinally often below the least stable ordinal (Section 1.4).
\item \textbf{Question.} Is $\Xp(\alpha) \le \Xo(\Xo(\alpha))$ for every infinite $\alpha$? Theorem 1 gives
  equality on $[\mu_0, \gamma)$.
\item \textbf{Question.} What is the least ordinal above $\gamma$ at which $\Xp \ne \Xo$? Lemma 7.3 does not
  apply to sets $Z$ that need a parameter from $L_u$ outside $H$.
\item \textbf{Question.} Theorem 3 (c) determines the least ordinal that shares all its $\Sigma^1_1$ properties
  with a larger one. What is the corresponding ordinal for other classes of properties, for instance the
  first-order properties of $L_{\alpha}$?
\end{enumerate}

\section{Verification status}

\begin{center}\begin{tabular}{p{0.24\linewidth}p{0.70\linewidth}}
\toprule
Part & Status \\
\midrule
All results & Written proofs, not checked by a person. AI systems checked the arguments;
that is not specialist review. The argument for $\delta_{\sigma} < \sigma^+$ in Proposition 6 was proposed by the AI
checker and has not been independently checked. \\
Formalisation & No end-to-end Lean formalisation. Two auxiliary facts are formalised in Lean: upward persistence of
$\Sigma_1$ formulas between transitive sets, used in Lemma 6.1, and $\dz{\alpha} \le \alpha^+$. \\
Imports & (L1)--(L4): Barwise and Devlin; the locators were not checked against the books. (L5), Theorem I.6
(used only for $\alpha < \sigma$), III.10, Remark IV.5 and the example after IV.6: Ressayre, read in a scanned copy
with gaps. Ressayre III.6 (a) is cited for comparison only. Girard's normal form and definition of a predilator: Girard (1981), Theorem 2.3.12 and
Definition 4.4.1. Admissibility of $\omega_1^y$: Sacks, no locator. \\
Other references & Hoole's paper is cited for its gap theorem only. Schlutzenberg's result is quoted from Schweber.
Of Aguilera's paper only Section 2 was read, and of Volume 2 of Girard's book only Annexes 11.B and 11.C;
Aguilera and Pakhomov and Girard and Ressayre were not read. \\
Novelty & Not established (Section 1.4). \\
\bottomrule
\end{tabular}\end{center}

No human specialist has reviewed the proofs. The passages a specialist is most likely to question are:
\begin{itemize}
\item the rank bound in Lemma 3.2 (ii) and its use in Theorem 2 (a);
\item the coding of $L_{\beta}$ in Lemma 3.3 (L);
\item Lemma 5.1 (c), where a model that is not well-founded computes a term correctly;
\item the diagonal sentence in Lemma 7.3, and the proof of Theorem 1 (i) for $r(d) < u$, which applies that lemma
  at $r(d)$ and uses that $Z_d$ is $\Sigma_1$-definable from $r(d)$;
\item Lemma 7.1, in particular the hypothesis ``countable in $L_{r^+}$'' as used in Step 3, and Step 1, where
  $F \circ F$ is considered on classes of infinite ordinals only;
\item the $\Sigma^1_1$ class used for $\Xp(\gamma) = \Xo(\gamma)$ in the proof of Theorem 1;
\item the elimination of parameters in the proof of Proposition 6.
\end{itemize}

\section*{Appendix A. Proof of Lemma 3.3}

\begin{proof}[Proof of Lemma 3.3] (C) The symbols are ternary $\mathrm{Pl}$, $\mathrm{Ti}$, $\mathrm{ent}$ and binary $\mathrm{len}$.
Let $N(x)$ say that every $y \le x$ is least or has an immediate predecessor; $N$ also denotes the set of these $x$, and let $x'$ be the successor of $x$.
$\theta_c$ is the conjunction of the following.
\begin{itemize}
\item[N1.] $N$ has no greatest element.
\item[N2.] $\mathrm{Pl}$ and $\mathrm{Ti}$ are graphs of binary functions on $N$ satisfying $x + 0 = x$,
  $x + y' = (x + y)'$, $x \cdot 0 = 0$, $x \cdot y' = x \cdot y + x$.
\item[N3.] $\mathrm{len}$ is a partial function into $N$; $\mathrm{ent}(s, i, a)$ implies $i < \mathrm{len}(s)$; and
  for each $i < \mathrm{len}(s)$ there is exactly one $a$ with $\mathrm{ent}(s, i, a)$, written $s_i$.
\item[N4.] Some $s$ has $\mathrm{len}(s) = 0$.
\item[N5.] For every $s$ in the domain of $\mathrm{len}$ and every $a$ there is $s^+$ with
  $\mathrm{len}(s^+) = \mathrm{len}(s)'$, $s^+_i = s_i$ for $i < \mathrm{len}(s)$, and last entry $a$.
\end{itemize}
In a well-order, $N$ is an initial segment whose nonzero elements are successors, so it has type $\le \omega$, and
type $\omega$ by N1. So a finite $A$ has no expansion. By N2 and induction $\mathrm{Pl}$ and $\mathrm{Ti}$ are the
true operations, so an arithmetical formula relativised to $N$ has its true meaning, and a recursive relation is
given by its $\Sigma_1$ definition. By N4, N5 and induction on the length, every finite sequence from $A$ is
$(s_0, \dots, s_{k-1})$ for some $s$ with $\mathrm{len}(s) = k$. Let $\operatorname{dom}_P(x)$ say that
$x_0 \in N$ is a trace, $\mathrm{len}(x) = 1 + \mathrm{ar}(x_0)$ and $x_1 < x_2 < \dots$; and let $x <_P y$ say
that both satisfy $\operatorname{dom}_P$ and the comparison of $P$ gives $<$ on $x_0$, $y_0$ and the order pattern
of the two tuples, which is coded by a number. The map is $x \mapsto (x_0, (x_1, x_2, \dots))$. Conversely, an
infinite well-order $A$ admits an injection of its finite sequences into $A$.

(L) Fix a primitive recursive coding of the formulas of $\{\in\}$ in variables $v_0, v_1, \dots$ in which the
index of every variable of a formula and the code of every subformula are at most the code of the formula. The further
symbols are binary $E$, $\mathrm{Lev}$ and ternary $\mathrm{Sat}$. Let $\bar 0$ be the least element of $A$ and let
$B_a(y)$ stand for $\exists a' < a\ \mathrm{Lev}(a', y)$. An \emph{$a$-assignment for $n$} is a code $s$ with
$\mathrm{len}(s) > n$ all of whose entries satisfy $B_a$; $s[i{:=}y]$ is $s$ with entry $i$ replaced by $y$, and a statement about $s[i{:=}y]$ is read as a statement about
every code of that sequence. Let
$\mathrm{Df}(a, n, s, y)$ stand for $B_a(y) \wedge \mathrm{Sat}(a, n, s[0{:=}y])$. $\theta_L$ is $\theta_c$ together with:
\begin{itemize}
\item[S1.] $E$ is extensional.
\item[S2.] $\mathrm{Lev}(\bar 0, x)$ if and only if $x$ has no $E$-predecessor, and such an $x$ exists.
\item[S3.] For $a > \bar 0$, a formula code $n \in N$ and an $a$-assignment $s$ for $n$: $\mathrm{Sat}(a, n, s)$
  obeys Tarski's clauses for the structure $(\{y : B_a(y)\}, E)$; in particular, for the code $n$ of
  $\exists v_i\, \varphi$, $\mathrm{Sat}(a, n, s)$ holds if and only if there is $y$ with $B_a(y)$ and
  $\mathrm{Sat}(a, \ulcorner\varphi\urcorner, s[i{:=}y])$.
\item[S4.] For $a > \bar 0$: $\mathrm{Lev}(a, x)$ if and only if there are $n$ and an $a$-assignment $s$ for $n$
  with $\forall y\, (y \E x \leftrightarrow \mathrm{Df}(a, n, s, y))$.
\item[S5.] For $a > \bar 0$ and $n$, $s$ as in S4 there is $x$ with $\forall y\, (y \E x \leftrightarrow \mathrm{Df}(a, n, s, y))$.
\item[S6.] Every $x$ satisfies $\mathrm{Lev}(a, x)$ for some $a$.
\end{itemize}
Let $A$ be a well-order with such an expansion, and write $\bar a$ for the order type of $\{a' : a' < a\}$. If
$\mathrm{Lev}(a, x)$ and $y \E x$ then $\mathrm{Lev}(a', y)$ for some $a' < a$, by S2 and S4. So the least $a$ with
$\mathrm{Lev}(a, x)$, which exists by S6, decreases along $E$. Hence $E$ is well-founded, each set
$B_a = \{y : B_a(y)\}$ is closed under $E$-predecessors, and by S1 the transitive collapse $\varpi$ of $(A, E)$ is
injective. We show $\varpi[B_a] = L_{\bar a}$ by induction on $a$; the induction hypothesis is
$\varpi[B_{a'}] = L_{\bar a'}$ for all $a' < a$. Let $a' < a$. If $a' = \bar 0$, then $\varpi$
maps $\{x : \mathrm{Lev}(\bar 0, x)\}$ onto $\{\emptyset\} = L_1$ by S2. If $a' > \bar 0$, induction on the formula
shows that $\mathrm{Sat}(a', n, s)$ holds if and only if $(B_{a'}, E)$ satisfies the formula with code $n$ under
$s$, by S3; here every $s[i{:=}y]$ has a code. So by S4 and S5, $\varpi$ maps $\{x : \mathrm{Lev}(a', x)\}$ onto the
set of definable subsets of $L_{\bar a'}$, which is $L_{\bar a' + 1}$. The union over $a' < a$ is $L_{\bar a}$. The
same union over all of $A$ gives $\varpi[A] = L_{\beta}$. Conversely, for infinite $\beta$ pull back $\in$, the
levels and satisfaction along a bijection between $\beta$ and $L_{\beta}$.

(K) The further symbols are $E$ and a binary $\mathrm{Sat}'$. $\theta_T$ is $\theta_c$, Tarski's clauses for
$\mathrm{Sat}'$ over the whole structure $(A, E)$, and the statement that $\mathrm{Sat}'(n, s)$ holds for every
code $n \in N$ of an axiom of $T$ and every sufficiently long code $s$. As in (L), $\mathrm{Sat}'$ is true
satisfaction on the codes in $N$, which are standard. The converse takes true satisfaction.
\end{proof}

\section*{Appendix B. Ressayre's class of functors}

\begin{statement}\head{Lemma B.1} For every recursive pattern set $\mathcal S$ there is a primitive recursive
pattern set $\mathcal S^c$ such that $\mathcal S^c(n)$ is finite for finite $n$, the functor
$n \mapsto P_{\mathcal S^c}(n)$ is primitive recursive on objects and maps, and
$\operatorname{rk} \mathcal S^c(b) = \operatorname{rk} \mathcal S(b)$ for every infinite well-order $b$.
\end{statement}

The tree $\mathcal S(n)$ can be infinite for finite $n$, because the depth of nodes and the data in a row can be
unbounded; a row without slots that may be repeated gives infinite depth already at $n = 0$. Also, membership in
$\mathcal S$ is only recursive. In $\mathcal S^c$ each row therefore carries a computation that verifies membership,
which makes the test primitive recursive,
and a block of \emph{extra slots}, at least as many as the size of what the row contains, whose labels must all be
different. A finite order runs out of labels, and this bounds every node.

\begin{proof} A row of a pattern in $\mathcal S^c$ is a row of a pattern in $\mathcal S$ whose tag is extended by
a halting computation showing that the pattern up to this row is in $\mathcal S$, followed by a block of extra
slots whose number exceeds the code of the extended tag and the number of original slots. It is required that the
labels in the extra slots of a node strictly increase in their order of occurrence through the whole node; the
order between extra labels and original labels is not restricted. On a finite order $n$ a node has at most $n$
extra labels, so the number of its rows and the codes of its tags are bounded by $n$, each row has at most $2n$ slots, and
membership is a bounded test.

Let $b$ be an infinite well-order. Deleting the computations and the extra labels maps each node of
$\mathcal S^c(b)$ to a node of $\mathcal S(b)$ and each child to a child. So by induction on the rank, the rank of
a node of $\mathcal S^c(b)$ is at most the rank of its image, and
$\operatorname{rk} \mathcal S^c(b) \le \operatorname{rk} \mathcal S(b)$; an infinite branch is mapped to an
infinite branch. Conversely, call a node of $\mathcal S^c(b)$ \emph{low} if its extra labels are among the first
$\omega$ elements of $b$. If $s^c$ is a low node with image $s$ and $s'$ is a child of $s$ in $\mathcal S(b)$, then
$s^c$ has a low child with image $s'$: extend the tag of the last row of $s'$ by the computation, which exists
because the pattern of $s'$ is in $\mathcal S$, and take the finitely many extra labels among the first $\omega$
elements of $b$ above the extra labels of $s^c$. The root is low. So by induction on the rank of $s$, a low node
with image $s$ has rank at least the rank of $s$, and an infinite branch of $\mathcal S(b)$ lifts to one of
$\mathcal S^c(b)$. Hence $\operatorname{rk} \mathcal S^c(b) \ge \operatorname{rk} \mathcal S(b)$.
\end{proof}

\begin{statement}\head{Lemma B.2} Let $F \in \mathcal R$ or $F \in \Pre$. There is $G \in \Pre$ that also belongs
to $\mathcal R$ such that for every infinite ordinal $b$: $G(b)$ is a well-order if and only if $F(b)$ is, and then
$|G(b)| > |F(b)|$.
\end{statement}

\begin{proof} A \emph{representative} over $b$ is a triple $(n, z, \vec a)$ with $z \in F(n)$, or $z$ a trace of
arity $n$ if $F \in \Pre$, and $\vec a$ an increasing $n$-tuple from $b$; it denotes the image of $z$ in $F(b)$
under the embedding $\vec a$. Every element of $F(b)$ has a representative, and two representatives are compared
inside $F(k)$ for $k$ the number of elements of $\vec a \cup \vec a'$, by a test that reads only the order pattern.
Let $\mathcal S$ consist of the patterns with tags $(n, z)$ and $n$ slots in each row such that the denoted
elements strictly descend. An infinite branch of $\mathcal S(b)$ is a descending sequence in $F(b)$, and if $F(b)$
is a well-order then by induction a node whose last row denotes an element of rank $\xi$ in $F(b)$ has rank at
least $\xi$. So $\operatorname{rk} \mathcal S(b) \ge |F(b)|$. Put $G = P_{\mathcal S^c}$ and apply Lemmas 3.1 and
B.1. $G(n)$ is a finite linear order, primitive recursively in $n$, and $G$ commutes with direct limits, so
$G \in \mathcal R$ after each $G(n)$ is identified with its order type.
\end{proof}

\section*{References}

\begin{itemize}
\item J. P. Aguilera, The order of reflection, arXiv:1906.11769 (Section 2).
\item J. P. Aguilera and F. Pakhomov, The $\Pi^1_2$ consequences of a theory, arXiv:2109.11652.
\item J. Barwise, \emph{Admissible Sets and Structures}, Springer, 1975.
\item J. Barwise, R. O. Gandy and Y. N. Moschovakis, The next admissible set, \emph{Journal of Symbolic Logic} 36 (1971), 108--120.
\item K. J. Devlin, \emph{Constructibility}, Springer, 1984.
\item J.-Y. Girard, $\Pi^1_2$-logic, Part 1: Dilators, \emph{Annals of Mathematical Logic} 21 (1981), 75--219.
\item J.-Y. Girard, \emph{Proof Theory and Logical Complexity}, Volume 1, Bibliopolis, 1987; Volume 2, unpublished
  manuscript, of which Part III (Chapters 8--12) is on the author's web page.
\item J.-Y. Girard and J.-P. Ressayre, \'El\'ements de logique $\Pi^1_n$, in \emph{Recursion Theory}, Proceedings
  of Symposia in Pure Mathematics 42, American Mathematical Society, 1985, 389--445.
\item F. G. Abramson and G. E. Sacks, Uncountable Gandy ordinals, \emph{Journal of the London Mathematical Society} (2) 14 (1976), 387--392.
\item R. O. Gandy, Proof of Mostowski's conjecture, \emph{Bulletin de l'Acad\'emie Polonaise des Sciences} 8 (1960), 571--575.
\item R. Gostanian, The next admissible ordinal, \emph{Annals of Mathematical Logic} 17 (1979), 171--203.
\item M. R. R. Hoole, The least $\alpha$ for which $E(\alpha)$ is inadmissible, \emph{Journal of the Australian
  Mathematical Society (Series A)} 41 (1986), 143--151.
\item J.-P. Ressayre, Bounding generalized recursive functions of ordinals by effective functors; a complement to
  the Girard theorem, in \emph{Proceedings of the Herbrand Symposium (Logic Colloquium '81)}, J. Stern (ed.),
  North-Holland, 1982, pp.\ 251 ff. (Theorem I.6; III.1--III.10; IV.0--IV.8.)
\item G. E. Sacks, \emph{Higher Recursion Theory}, Springer, 1990.
\item F. Schlutzenberg, answer 381041 on MathOverflow, 2021.
\item N. Schweber, Clarifying ordinals, arXiv:2408.10367 (Section 3.2).
\end{itemize}

\end{document}
